% Created 2026-08-25 Tue 08:10
% Intended LaTeX compiler: lualatex
\documentclass[11pt]{article}

\usepackage{amsmath}
\usepackage{fontspec}
\usepackage{graphicx}
\usepackage{longtable}
\usepackage{wrapfig}
\usepackage{rotating}
\usepackage[normalem]{ulem}
\usepackage{capt-of}
\usepackage{hyperref}
\usepackage{minted}
\def\EDITMODE{}
\input{~/org/templates/org-export-general-template.tex}
\loadEquity
\setcounter{secnumdepth}{5}
\author{Andrew Jensen}
\date{2026-08-21 2026-08-24}
\title{Rust direct-rs ↔ C++ FFI Handshake (reference)}
\hypersetup{
 pdfauthor={Andrew Jensen},
 pdftitle={Rust direct-rs ↔ C++ FFI Handshake (reference)},
 pdfkeywords={},
 pdfsubject={},
 pdfcreator={},
 pdflang={English}}
\usepackage{biblatex}
\addbibresource{~/org/biblio.bib}
\begin{document}

\maketitle

\section{Purpose}
\label{sec:orgf21adce}

Reference for replacing the C++ DIRECT core (\texttt{DirectOptimizer}) with the Rust
implementation in \texttt{rust/direct-rs}, keeping the C++ coordinator
(\texttt{OptimizerManager}, \texttt{CostFunctionManager}, GPU cost) intact. This is a
\textbf{handshake spec} — the concrete FFI boundary Rust and C++ must agree on — not a
full implementation plan. It records every option considered so the reader can
choose deliberately.
\subsection{How to use this doc (learning path)}
\label{sec:orgb153688}

This is a living handshake \textbf{and} a tutor. Every section that teaches a Rust
concept is tagged with three pointers:

\begin{itemize}
\item \textbf{\textbf{Brown Book}} — the free, interactive, visual book you prefer, at
\url{https://rust-book.cs.brown.edu/}. A section like ch 04-02 opens at
\url{https://rust-book.cs.brown.edu/ch04-02-references-and-borrowing.html}.
The Brown variant adds quizzes and explainers the static book lacks.
\item \textbf{\textbf{rust-skills}} — one of the 265 prioritized idioms at
\url{../../.pi/agent/skills/rust-skills/SKILL.md} (rule ids like
\texttt{own-borrow-over-clone}, one .md per rule).
\item \textbf{\textbf{normie ↔ idiomatic}} — the shape a C++/Python developer reaches for first,
and the idiomatic Rust you actually want.
\end{itemize}

The through-line question for every block: \textbf{``what did the Rust type system
just make easier for me, and what did it make impossible?''} If a section
reads as ``translate C++ into Rust'', that's the normie path — re-read it for
the place where ownership/borrowing, or exhaustive =match=/enums, or
iterator pipelines change what ``good'' means.

Carry the core line from the Brown book (its ch 4):

> Rust's central feature is \textbf{ownership}, and it pervades even tiny examples.

Coming from \textbf{\textbf{Python}}: ch 3-8 is mostly review (variables, functions,
structs, control flow). Camp on ch 4 (ownership), ch 9 (errors),
ch 10 (traits/generics), ch 13 (closures/iterators), ch 15 (smart pointers),
ch 17 (trait objects), ch 19 (unsafe).

Coming from \textbf{\textbf{C++}}: you'll feel at home in ch 15 (\texttt{Box} ≈ \texttt{unique\_ptr}) and
ch 19 (\texttt{unsafe} ≈ raw pointers), and least at home in ch 4 (borrowing) —
Rust's borrow checker is stricter than =const\&- discipline.

Suggested first pass (each 10-25 min, do the quizzes):

\begin{enumerate}
\item ch 04 — ownership/borrow/slices; the whole doc adheres to this.
\item ch 09-02 — \texttt{Result} + \texttt{?}; the error-channel section arises from it.
\item ch 10-02 — traits; the \texttt{Cost} trait and \texttt{IOptimizer} options build on it.
\item ch 13 — closures/iterators; batching and POH gathering lean on it.
\item ch 15-01 — \texttt{Box}; where \texttt{rust::Box<DirectOptimizer>} comes from.
\item ch 19-01 — \texttt{unsafe}; what \texttt{unsafe extern} actually gives up.
\end{enumerate}
\section{Toolchain first — make the compiler your tutor}
\label{sec:org55a873f}

You cannot learn Rust by reading. The borrow checker is the teacher and it only
speaks when you compile. Before Milestone 0, build a loop where you can ask the
compiler a question and get an answer in under three seconds. Everything below
this section assumes that loop exists.
\subsection{The crate layout this doc implies}
\label{sec:org8e4a061}

\begin{minted}[]{text}
rust/direct-rs/
  Cargo.toml
  build.rs              # cxx-build: compiles the generated C++ glue
  src/
    lib.rs              # #[cxx::bridge] ONLY — the boundary, nothing else
    pose.rs             # Pose, Direction                      (Steps 1-2)
    hyperbox.rs         # Hyperbox, size(), trisect             (Steps 3, 7)
    tree.rs             # BoxTree: two-level BTreeMap, gather, reinsert (Steps 4-5)
    poh.rs              # convex_hull — pure                   (Step 6)
    cost.rs             # trait Cost + AnalyticCost stub + CppCost impl (M6)
    error.rs            # OptimizerError (thiserror)           (M8)
    optimizer.rs        # DirectOptimizer::run — the loop      (M9)
  tests/
    parity.rs           # Tier-1 analytic golden vs C++        (Appendix F)
\end{minted}

Why split it this way: \texttt{lib.rs} holding \textbf{only} the bridge is FFI containment —
the unsafe surface is one file you can audit. Every other module is ordinary
safe Rust, unit-testable with no C++ and no GPU in the room. \textbf{\textbf{If a module
can't be tested headlessly, FFI has leaked into it.}} That single rule is worth
more than any diagram in this doc.
\subsection{Cargo.toml}
\label{sec:org0aa3014}

\begin{minted}[]{toml}
[package]
name = "direct-rs"
version = "0.1.0"
edition = "2021"

[lib]
crate-type = ["staticlib", "rlib"]   # staticlib for CMake to link; rlib for Rust tests

[dependencies]
cxx = "1.0"
ordered_float = "4"
thiserror = "2"

[build-dependencies]
cxx-build = "1.0"

[dev-dependencies]
approx = "0.5"      # assert_relative_eq! for float parity
proptest = "1"      # property tests for hull + trisect invariants
\end{minted}

Note \texttt{[dev-dependencies]}: these compile for \texttt{cargo test} only and never ship
into the C++ binary. That separation is why a test stub cost (\texttt{AnalyticCost})
costs the production build nothing.
\subsection{build.rs (the other half of the handshake)}
\label{sec:orgf0bdaf7}

\begin{minted}[]{rust}
fn main() {
    cxx_build::bridge("src/lib.rs")
        .include("../../include")            // where domain/cost.h lives
        .flag_if_supported("-std=c++17")
        .compile("direct_rs");
    println!("cargo:rerun-if-changed=src/lib.rs");
    println!("cargo:rerun-if-changed=../../include/domain/cost.h");
}
\end{minted}

The header C++ will \texttt{\#include} lands at:

\begin{minted}[]{text}
target/<profile>/cxxbridge/direct-rs/src/lib.rs.h
\end{minted}

Open that file on day one. It is the ground truth for every ``verify the
verbatim signature'' warning in this doc — \texttt{rust::Box<T>}, \texttt{rust::Slice<const
T>}, the generated \texttt{Pose} struct, all of it. You never have to guess.
\subsection{The six commands that are your feedback loop}
\label{sec:org4a84b2f}

\begin{center}
\begin{tabular}{ll}
Command & When / why\\
\hline
\texttt{cargo check} & every 30 seconds — type-checks without codegen; fastest tutor\\
\texttt{cargo clippy -{}-{} -W clippy::pedantic} & after it compiles — clippy \textbf{names} the idiom you missed\\
\texttt{cargo test} & after each Step; \texttt{cargo test hull} filters by name\\
\texttt{cargo expand -{}-{}lib ffi} & to see what \texttt{\#[cxx::bridge]} actually generated\\
\texttt{cargo doc -{}-{}open} & read \texttt{ordered\_float} / \texttt{cxx} docs offline, versioned to you\\
\texttt{cargo fmt} & never argue about layout; the formatter is canonical\\
\end{tabular}
\end{center}

Also bind: \texttt{rustc -{}-{}explain E0502} (full prose for any error code) and
\texttt{rustup doc -{}-{}book} (the offline book, when you want the non-Brown text).

A note on clippy for a learner: run it with \texttt{-W clippy::pedantic} even though
it's noisy. Half of the 265 rust-skills rules are things clippy will tell you
for free, with a link. Treat every clippy lint as a flashcard: read the name
(\texttt{needless\_range\_loop} → \texttt{perf-iter-over-index}), fix it, move on.
\subsection{Read the error, not the code}
\label{sec:orgd2631f0}

A rustc error has four parts: the code (\texttt{E0502}), the one-line claim, the
annotated span, and a \texttt{help:} line that is usually the literal fix. Beginners
stare at the span. The fastest learners read the \texttt{help:} first, then
\texttt{rustc -{}-{}explain}, and only then look at their own code. Appendix A's table and
the glossary's error lookup are the project-specific version of that habit.
\section{Learning Roadmap (plan it, then fill it)}
\label{sec:orgf9a092b}

This doc is big enough that the order you read it \textbf{is} a learning path. Use
this roadmap as your map; the section headings are the terrain. Each milestone
ends with a self-check you should be able to answer without looking before
moving on.
\subsection{Milestone 0 — why this is a two-language project}
\label{sec:org751644c}

Goal: before any Rust, understand the split. The C++ side owns the cost and
lifecycle; Rust owns the search algorithm. Everything else is plumbing between
two runtimes.

\begin{itemize}
\item Read: \texttt{Purpose}, \texttt{Contract being filled}, \texttt{One-picture flow}.
\item \textbf{\textbf{Book}}: none — orientation. Return to it after every milestone; the
handshake diagram is the anchor.
\item Self-check: state which side owns (1) the cost function, (2) the optimizer's
search loop, (3) the worker thread, (4) the UI progress, and (5) the
cumulative budget.
\end{itemize}
\subsection{Milestone 1 — value types: Pose + Hyperbox}
\label{sec:org516a1f0}

Goal: model the data before any algorithm. Internalize \texttt{Copy} vs \texttt{Clone} vs
move, and why small POD structs cross FFI by value.

\begin{itemize}
\item Concept: \texttt{\#[derive(Copy, Clone)]} on small structs; what a value is in Rust.
\item \textbf{\textbf{Book}}: ch 04-01 (moves), ch 05-01 (structs), ch 06 (enums for \texttt{Direction}).
\item \textbf{\textbf{rust-skills}}: \texttt{own-copy-small}, \texttt{own-move-large}, \texttt{type-enum-states}.
\item Self-check: explain why \texttt{Pose} is \texttt{Copy} but a \texttt{Hyperbox} whose center is a
\texttt{Pose} stays \texttt{Copy} (small POD), vs one that wraps a \texttt{Vec} which isn't — i.e.
what makes a type \texttt{Copy}-eligible and what \texttt{Copy} costs you.
\end{itemize}
\subsection{Milestone 2 — the BTreeMap design}
\label{sec:org3fe326b}

Goal: the central data-structure decision, and the most Rust-flavored one in
the doc. Learn \texttt{Ord}-based containers and the \texttt{OrderedFloat} wrapper.

\begin{itemize}
\item Concept: \texttt{f64} is not \texttt{Ord}; \texttt{BTreeMap} needs a total order; wrapping floats.
\item \textbf{\textbf{Book}}: ch 08 (collections) + the \texttt{ordered\_float} crate docs (the Brown
book has no \texttt{OrderedFloat} section — this is where you learn to read a
dependency's docs).
\item \textbf{\textbf{rust-skills}}: \texttt{coll-map-choice}, \texttt{coll-set-membership}, \texttt{num-float-compare}.
\item Self-check: why can't you put \texttt{f64} directly as a \texttt{BTreeMap} key, and what is
the minimum wrapper that fixes it?
\end{itemize}
\subsection{Milestone 3 — the gather (read-only borrow)}
\label{sec:orgc126bce}

The first actual pass over the tree. Internalize \textbf{read-only borrow} = a \texttt{\&self}
pass with no mutation and no clone.

\begin{itemize}
\item Concept: \texttt{iter()}, \texttt{first\_key\_value()}, projecting \texttt{(size, cost)} without
touching the box.
\item \textbf{\textbf{Book}}: ch 04-02 (borrowing), ch 13-02 (iterator adapters).
\item \textbf{\textbf{rust-skills}}: \texttt{own-borrow-over-clone}, \texttt{perf-iter-over-index}.
\item Self-check: why does the gather take \texttt{\&self} (not \texttt{\&mut}), and what would
break in the type system if you tried to \texttt{.clone()} the box out?
\end{itemize}
\subsection{Milestone 4 — POH selection (pure transform)}
\label{sec:org4cbe915}

Goal: \texttt{convex\_hull(\&[PohPoint]) -> Vec<PohPoint>} as a pure function. Internalize
value-in/value-out and the absence of side effects.

\begin{itemize}
\item Concept: mapping (size, cost) pairs through a hull algorithm; returning owned
\texttt{Vec<PohPoint>}.
\item \textbf{\textbf{Book}}: ch 10-02 (traits) + ch 13 (closures as data).
\item \textbf{\textbf{rust-skills}}: \texttt{anti-vec-for-slice}, \texttt{perf-collect-once}.
\item Self-check: if \texttt{convex\_hull} is pure, what do you pass in and what comes out,
and how do you know the tree wasn't touched?
\end{itemize}
\subsection{Milestone 5 — trisection + the ownership dance}
\label{sec:orgf7e9ebd}

The crown jewel of the doc and the deepest borrowing lesson: take = move,
trisect = owned transform, reinsert = move back.

\begin{itemize}
\item Concept: \texttt{pop\_first()} returns owned; \texttt{\&mut self} is exclusive.
\item \textbf{\textbf{Book}}: ch 04 (ownership) — the payoff milestone; re-read it.
\item \textbf{\textbf{rust-skills}}: \texttt{own-move-large}, \texttt{mem-take-replace}, \texttt{own-arc-shared}.
\item Self-check: why must you pop/remove the winner OUT of the map before
trisecting it, and why can't you mutate a map-resident box in place?
\end{itemize}
\subsection{Milestone 6 — cost injection + the trait}
\label{sec:orgaa54efa}

Goal: the \texttt{Cost} capability; batch-vs-serial; the one-FFI-hop idea.

\begin{itemize}
\item Concept: \texttt{trait Cost}, \texttt{impl Fn} vs \texttt{Box<dyn Cost>}.
\item \textbf{\textbf{Book}}: ch 10-02 (traits), ch 17 (trait objects), ch 13 (closures).
\item \textbf{\textbf{rust-skills}}: \texttt{trait-dyn-vs-generic}, \texttt{closure-fn-trait-bounds}.
\item Self-check: why is \texttt{Cost} a trait and not a concrete struct field, and when
would you prefer \texttt{impl Fn} vs \texttt{Box<dyn Cost>}?
\end{itemize}
\subsection{Milestone 7 — the FFI boundary (cxx)}
\label{sec:org20d1bae}

Goal: see your Rust reach into C++. Own \texttt{\#[cxx::bridge]}, opaque types, shared
structs/enums, the no-panic rule.

\begin{itemize}
\item Concept: \texttt{extern "Rust"} vs \texttt{extern "C++"}, \texttt{Box<DirectOptimizer>} on C++ side,
shared enum \texttt{> wildcard =\_} arm.
\item \textbf{\textbf{Book}}: ch 19 (unsafe/FFI), ch 16 (threads for Send/Sync).
\item \textbf{\textbf{rust-skills}}: \texttt{unsafe-safety-comment}, \texttt{err-no-unwrap-prod}.
\item Self-check: what two things can never cross a cxx boundary, and why does that
force the status-struct + ProgressSink shapes?
\end{itemize}
\subsection{Milestone 8 — errors + status (the gate)}
\label{sec:org31aa25c}

Goal: the error story inside (\texttt{Result} + =?) vs at the edge (status struct).

\begin{itemize}
\item Concept: two-tier errors; Result in, status out; panic never crosses.
\item \textbf{\textbf{Book}}: ch 09 (all of it).
\item \textbf{\textbf{rust-skills}}: \texttt{err-result-over-panic}, \texttt{err-thiserror-lib}, \texttt{err-question-mark}.
\item Self-check: converting \texttt{OptimizerError} to a \texttt{RunOutcome.ok} — what is lost
and what is kept, and why is that acceptable at the FFI edge?
\end{itemize}
\subsection{Milestone 9 — assemble the loop + A/B}
\label{sec:orgef685ed}

The finish: wire gather → POH → trisect → reinsert into one \texttt{run()} behind the
\texttt{OptimizerBackend} A/B switch.

\begin{itemize}
\item Concept: glue, ownership flow end-to-end, the \texttt{ok} gate.
\item \textbf{\textbf{Book}}: ch 04 + ch 13 once more (full recap).
\item Self-check: walk the ownership of ONE hyperbox from insertion to trisection
to reinsertion, naming each borrow/move.
\end{itemize}

Return here after every milestone and re-check the handshake diagram: each one
fills a corner of it. The roadmap is the planning; the sections that follow are
the terrain.
\subsection{Milestone exercises — write these, don't read them}
\label{sec:orga29badd}

Each milestone above has a self-check \textbf{question}; here is the matching
\textbf{keystroke}. Do the exercise before you read the corresponding Step, then read
the Step and see what you'd have done differently. Nothing here needs C++, a
GPU, or the real cost function — all of it is \texttt{cargo test}-able in the crate.

\begin{center}
\begin{tabular}{lll}
\# & Exercise (write it, compile it, test it) & Done when\\
\hline
E1 & \texttt{Pose} + \texttt{Direction}; a \texttt{const DIRECTIONS: [Direction; 6]}; \texttt{Pose::shifted(dir, d)} & a test asserts a shift on X leaves Y..ZA untouched\\
E2 & \texttt{Hyperbox::size()}; unit-test that one trisection on the min axis shrinks size & \texttt{assert!(child.size() < parent.size())}\\
E3 & Deliberately try \texttt{BTreeMap<f64, Hyperbox>}; read the error; then fix with \texttt{OrderedFloat} & you can quote the trait that's missing (\texttt{Ord})\\
E4 & \texttt{BoxTree::insert} + \texttt{gather\_poh}; build a 5-box tree by hand in a test & gather returns one point per size-row, min cost\\
E5 & Try to \texttt{.clone()} a box out of \texttt{gather\_poh} and make it compile with \texttt{\&self} & you can explain \textbf{why} it forces \texttt{\&mut self}\\
E6 & \texttt{convex\_hull} on a hand-made 6-point set with a known hull & test passes AND the fn takes \texttt{\&[PohPoint]}\\
E7 & \texttt{trisect(parent) -> (Hyperbox, [Hyperbox; 2])}; assert depth sum increased by 1 & three children, parent unusable after (E0382 if you try)\\
E8 & \texttt{trait Cost} + \texttt{AnalyticCost} (e.g. sphere function \texttt{sum(x\_i\textasciicircum{}2)}) & \texttt{run()} finds the origin to 1e-6 with no C++\\
E9 & \texttt{OptimizerError} with \texttt{thiserror}; make \texttt{run} return \texttt{Result}, then flatten to \texttt{RunOutcome} & \texttt{ok=false} for a zero-range input, no panic\\
E10 & Wire the loop; run against \texttt{AnalyticCost} under a 20k budget & \texttt{calls + offset < budget} holds exactly, no overrun\\
\end{tabular}
\end{center}

E3 and E5 are the two that \textbf{look} like a waste of time and aren't: they are
deliberate failures. The point of writing code that doesn't compile is that the
error message is the lesson, and you will meet both errors again for real later
(see Appendix A). Budget \textasciitilde{}20 minutes each; the rest are 30-60.

A good rhythm: exercise → the matching Step in ``Designing DIRECT from scratch''
→ the Brown Book section it cites → \texttt{cargo clippy} on what you wrote. Four
passes over the same idea, each from a different angle.
\section{Contract being filled}
\label{sec:org144ee96}

The C++ \texttt{DirectOptimizer::Run()} contract, from \texttt{include/domain/direct\_optimizer.h}:

\begin{center}
\begin{tabular}{lll}
Method & Signature & Semantics\\
\hline
\texttt{Run()} & \texttt{bool Run()} & false=error; loop guard \texttt{(calls+offset)<budget \&\& !stop}\\
\texttt{GetCostFunctionCalls()} & \texttt{unsigned int} & \texttt{cost\_function\_calls\_ + call\_offset\_}\\
\texttt{GetNonFiniteCount()} & \texttt{unsigned int} & count of non-finite evals\\
\texttt{GetOptimumLocation()/Value()} & \texttt{Point6D / double} & denormalized physical best / best cost\\
\texttt{Stop()} & \texttt{void} & cooperative, checked between iterations\\
\texttt{SetCallOffset(u)} & \texttt{void} & cumulative offset; canonical caps trunk 20k → 2x branch 25k/30k → leaf 35k\\
\texttt{SetIterationCallback(f)} & \texttt{void} & after each ConvexHull+Trisect, \textasciitilde{}30fps\\
\texttt{SetImprovementCallback(f)} & \texttt{void} & on new best (emits UpdateOptimum)\\
\texttt{SetBatchCost(f)} & \texttt{void} & optional U11 batch sibling\\
\end{tabular}
\end{center}

\texttt{Options} struct (8 fields, all fail-fast unless default):
\texttt{selection} (SelectionMode=Original), =epsilon=0.0, =delta\textsubscript{limit}=bool(false),
=size\textsubscript{measure}=SizeMeasure=L2, =split\textsubscript{rule}=SplitRule=OneSide,
=ties=TieSelection=All, =hidden\textsubscript{constraints}=bool(false),
=globally\textsubscript{biased}=bool(false).

Truths that must survive any backend:

\begin{itemize}
\item cumulative \texttt{calls+offset} budget semantics (canonical caps: trunk 20k → 2x branch 25k/30k → leaf 35k, per direct\textsubscript{optimizer.h})
\item the SAME optimum for the same cost sequence (Tier-1 analytic golden is the judge)
\item eval ordering: calls++ → non-finite → optimum update → storage → improvement callback
\item \texttt{evaluate} executed \textbf{\textbf{serially}} on one worker thread (GPUModel pose is mutated per call)
\end{itemize}
\subsubsection{Rust learning — ``options as data'' vs the god object}
\label{sec:org18e2fee}

The C++ \texttt{Options} struct is already \textbf{value-shaped} (a struct passed by value)
— a good instinct. The normie-from-C++ error is to treat it as a bag of
bools/enums read by \texttt{if} inside \texttt{Run()}. Rust nudges you to make the impossible
state \textbf{unrepresentable} rather than merely checked.

\begin{itemize}
\item \textbf{\textbf{Book}}: ch 05 (structs), ch 06 (enums + exhaustive \texttt{match}),
ch 10-02 (traits).
\item \textbf{\textbf{rust-skills}}: \texttt{type-enum-states}, \texttt{type-newtype-validated},
\texttt{api-parse-dont-validate}, \texttt{api-common-traits}, \texttt{num-float-compare},
\texttt{anti-over-abstraction}.
\end{itemize}

Normie (a loose struct with raw codes; callers must remember valid values):

\begin{minted}[]{rust}
pub struct Options { pub selection: u8, pub epsilon: f64, pub delta_limit: bool }
// a caller writing selection: 8 compiles fine and is wrong at runtime
\end{minted}

Idiomatic (a named enum + derive so invalid values can't exist):

\begin{minted}[]{rust}
#[derive(Debug, Clone, Copy, PartialEq, Eq)]
pub enum SelectionMode { Original }      // only one legal value today
#[derive(Debug, Clone, Copy, Default)]
pub struct Options { pub selection: SelectionMode }  // Default == the C++ default
\end{minted}

Why \texttt{epsilon}, \texttt{size\_measure}, \texttt{split\_rule}, \texttt{ties}, \texttt{hidden\_constraints},
\texttt{globally\_biased} are \textbf{absent} above: the C++ side fails-fast on any
non-default value. The idiomatic Rust is to \textbf{\textbf{not declare the knob until you
implement it}} (\texttt{anti-over-abstraction}). Declaring a field you cannot honor
is a promise the type system can't keep; dropping it entirely is a
\texttt{compile-time} guarantee that beats the run-time \texttt{ValidateOptions()}.

The cumulative budget caps (trunk 20k → 2x branch 25k/30k → leaf 35k) are
better as \textbf{\textbf{data}} than scattered literals — a single \texttt{Budget} newtype
(\texttt{type-newtype-ids}) and the ladder defined once, used everywhere. If you ever
\textbf{compute} a budget, reach for \texttt{num-overflow-explicit} (checked/saturating
arithmetic) rather than letting a += overflow wrap.
\section{Designing DIRECT from scratch, step by step}
\label{sec:org9315bce}

This is the procedural heart. Forget the FFI and options for a moment: here we
derive the Rust DIRECT data model and loop \textbf{one step at a time}, in the order
you'd discover them while porting. Each step introduces exactly one idea,
usually one type, and checks itself before moving on. This is where the
``normie → idiomatic'' comparisons get their teeth — each step is a real decision,
not a style preference.

Throughout, ``deduce'' means \textbf{let the compiler + the borrow checker teach you},
not just copy a pattern. The Brown Book's ch 4 is the reference for steps 3-5;
re-open it when a step feels abstract.
\subsection{Step 1 — model a point: `Pose` (value type)}
\label{sec:orgd39e1de}

The whole algorithm is a search over 6-D poses. First question: what is a pose
\textbf{on its own}? A bundle of six \texttt{f64} whose identity is its contents. That is a
value type — and in Rust the first reflex is to ask: should this be \texttt{Copy}?

Normie-from-C++ (a class with methods + hidden members):
\begin{minted}[]{cpp}
class Pose { double x,y,z,xa,ya,za; /* getters, setters... */ };
\end{minted}

Idiomatic (a plain value struct, \texttt{Copy}, fields public):
\begin{minted}[]{rust}
#[derive(Clone, Copy, PartialEq, PartialOrd)]
pub struct Pose { pub x:f64, pub y:f64, pub z:f64, pub xa:f64, pub ya:f64, pub za:f64 }
\end{minted}

Why: \texttt{Pose} is 6×=f64= = 48 bytes — genuinely small — and its equality/order are
meaningful (\texttt{PartialEq}, \texttt{PartialOrd}). \texttt{Copy} is right for a small plain-data
value (rust-skills \texttt{own-copy-small}): passing it around copies 48 bytes (cheap),
and \texttt{Copy} lets you read it out of a map key without \texttt{.clone()} ceremony.

\begin{itemize}
\item \textbf{\textbf{rust-skills}}: \texttt{own-copy-small}, \texttt{type-enum-states}.
\item \textbf{\textbf{Book}}: ch 04-01 (Copy vs move), ch 05-01 (structs).
\end{itemize}

Stop. \textbf{\textbf{Check}} — is a \texttt{Pose} a thing you own or a thing you copy? It is both —
\texttt{Copy} is the honest label for ``duplicating it is cheap and preserves meaning,''
which tells the compiler it may copy at will.
\subsection{Step 2 — model an axis: `Direction` (enum, not ints)}
\label{sec:orgee4b68a}

DIRECT trisects along one of six axes. The normie input is ints 0..5 or named
constants; Rust has a better tool: \texttt{Direction} as an \texttt{enum}.

\begin{minted}[]{rust}
#[derive(Clone, Copy)] // tiny tag — keep it Copy
pub enum Direction { X_DIR, Y_DIR, Z_DIR, XA_DIR, YA_DIR, ZA_DIR }
\end{minted}

An \texttt{enum} buys three things the normie approaches lack:
\begin{enumerate}
\item no invalid value (no ``axis 7'') — \texttt{type-enum-states};
\item exhaustive \texttt{match} — the compiler forces all six cases; and
\item it pairs with an index into \texttt{depths: [u32; 6]} via a small const mapping
array instead of a cast. You already have `const DIRECTIONS: [Direction; 6]`.

\item \textbf{\textbf{rust-skills}}: \texttt{type-enum-states}, \texttt{pat-exhaustive-enum}.
\item \textbf{\textbf{Book}}: ch 06 (enums + match).
\end{enumerate}
\subsection{Step 3 — model a box: `Hyperbox`}
\label{sec:org189a25f}

Now the atom of the search. A hyperbox has a center (\texttt{Pose}), a per-axis depth
\texttt{depths: [u32; 6]} (how many times each axis has been trisected), and a cost at
that center (\texttt{Option<f64>} — we'll see why).

Normie-from-C++ (a class with a \texttt{size()} method + friends):
\begin{minted}[]{cpp}
class HyperBox6D { /* ... */ double GetSides(); /* ... */ };
\end{minted}

Idiomatic (a plain struct + a method on an \texttt{impl}):
\begin{minted}[]{rust}
#[derive(Clone, Copy)]
pub struct Hyperbox {
    pub cost_at_center: Option<f64>,
    pub center: Pose,
    pub depths: [u32; 6],
}
impl Hyperbox {
    /// L2 norm of box extent from depths — the "size" used by POH.
    pub fn size(&self) -> f64 {
        self.depths.iter()
            .map(|e| 3f64.powf(-2.0 * (*e as f64)))
            .sum::<f64>()
            .sqrt()
    }
}
\end{minted}

Two Rust habits to internalize here:

\begin{itemize}
\item \textbf{\textbf{\texttt{Option<f64>} for ``not scored yet.''}} The center box inherits its parent's
cost; the two shifted boxes don't exist yet — pending a batch eval. \texttt{Option}
\end{itemize}
makes ``not yet known'' an explicit state instead of a magic \texttt{DBL\_MAX=/}-1=
sentinel. This is \texttt{type-option-nullable} + \texttt{api-parse-dont-validate}: represent
  the absence; don't smuggle a fake number.
\begin{itemize}
\item \textbf{\textbf{A method computes what isn't stored.}} \texttt{size()} is derived from \texttt{depths}
\end{itemize}
every time; you don't cache it in the struct. Fewer fields = fewer places to be
inconsistent. (Add a cache only if it shows up in the profiler =
\texttt{perf-profile-first}.)

\begin{itemize}
\item \textbf{\textbf{rust-skills}}: \texttt{type-option-nullable}, \texttt{api-parse-dont-validate},
\texttt{own-copy-small}, \texttt{perf-profile-first}.
\item \textbf{\textbf{Book}}: ch 05-01 (structs), ch 06 (Option), ch 13-02 (iterators — the \texttt{sum}
is already an iterator pipeline).
\end{itemize}

Stop. \textbf{\textbf{Check}} — why \texttt{Option<f64>} for cost rather than \texttt{f64} with \texttt{-1}
``meaning missing''? Because \texttt{-1} might be a valid cost that gets silently used;
\texttt{None} is a kind the compiler forces you to handle before you can use the box.
\subsection{Step 4 — the two-level `BTreeMap` (the Ord problem)}
\label{sec:org12c6f45}

The heart of the data model is a container of hyperboxes keyed by size (outer)
then cost (inner). The moment you reach for a sorted map you hit a very
Rust-specific wall: \textbf{\textbf{\texttt{f64} is not \texttt{Ord}.}}

Why: \texttt{BTreeMap} (and \texttt{BTreeSet}, \texttt{BinaryHeap}, sort) needs a \textbf{total order} — the
\texttt{Ord} trait. But \texttt{f64} has \texttt{NaN}, which compares ``not equal to everything,
including itself,'' breaking total-order. So the standard library deliberately
refuses \texttt{f64} as a map key / in \texttt{Ord} contexts.

Normie-from-C++ (map on \texttt{std::map<double,...>} — C++ lets you, NaN foot-gun
later):
\begin{minted}[]{cpp}
std::map<double, std::vector<Box*>> columns;  // NaN quietly poisons comparisons
\end{minted}

Idiomatic (wrap the float + add an id for ties):
\begin{minted}[]{rust}
use ordered_float::OrderedFloat;             // newtype: f64 that IS Ord
pub type SizeKey = OrderedFloat<f64>;
pub type CostKey = (OrderedFloat<f64>, u64); // (cost, unique id) so equal costs don't collide
\end{minted}

\texttt{OrderedFloat<f64>} is a newtype that implements \texttt{Ord} by asserting the value is
never \texttt{NaN}. It's the idiomatic escape hatch when you know your floats are total
(your sizes/costs are never NaN). The \texttt{u64} tiebreak makes the inner key injective.

\begin{itemize}
\item \textbf{\textbf{rust-skills}}: \texttt{num-float-compare}, \texttt{type-newtype-ids}, \texttt{coll-map-choice}.
\item \textbf{\textbf{Book}}: the \texttt{ordered\_float} crate isn't in the Brown book — this is your
first ``read a dependency's docs'' moment. The concepts (traits, newtype, \texttt{Ord})
are ch 10-02 + ch 08.
\end{itemize}

Stop. \textbf{\textbf{Check}} — why is the inner key \texttt{(cost, u64)} rather than just \texttt{cost}?
Because two boxes can have the same cost; a \texttt{BTreeMap} keyed by cost alone
would let the second silently overwrite the first. The monotonically increasing
\texttt{u64} box id breaks the tie so both survive — \texttt{type-newtype-ids} in a container.
\subsection{Step 5 — the gather (read-only pass)}
\label{sec:orga685729}

Now the first read over the tree: collect, for every size-row, the min-cost
box's (size, cost) pair — the input to POH. This is a \textbf{\textbf{read-only borrow}} pass
= \texttt{\&self}, no mutation, no clones.

Normie-from-C++ (a \texttt{const\&} method that iterates + copies minima out):
\begin{minted}[]{cpp}
std::vector<std::pair<double,double>> gather(const std::map<...>& rows) {...}
\end{minted}

Idiomatic (read-only \texttt{iter()} + \texttt{first\_key\_value()}, projecting two numbers):
\begin{minted}[]{rust}
fn gather_poh(&self) -> Vec<PohPoint> {
    let mut rows = Vec::new();
    for (size_key, cost_val) in self.boxes.iter() {          // &, immutable
        if let Some((key, _box)) = cost_val.first_key_value() {  // &, immutable
            rows.push(PohPoint {
                size: (*size_key).into_inner(),   // copy the f64 out of the key
                cost: (key.0).into_inner(),       // copy the cost, not the box
            });
        }
    }
    rows
}
\end{minted}

The teaching gold is \textbf{what did NOT happen}: we never \texttt{.clone()} a \texttt{Hyperbox},
we never mutated the tree, so \texttt{gather\_poh} is \texttt{\&self} and it's callable while
other borrows are live. If the normie version cloned boxes into a
\texttt{Vec<Hyperbox>}, it'd be forced into \texttt{\&mut self} (or fight the borrow) and carry
the whole box when only two numbers were needed. The borrow checker \textbf{points to}
the leaner shape.

\begin{itemize}
\item \textbf{\textbf{rust-skills}}: \texttt{own-borrow-over-clone}, \texttt{perf-iter-over-index},
\texttt{own-slice-over-vec}.
\item \textbf{\textbf{Book}}: ch 04-02 (borrowing — the whole step), ch 13-02 (iterator
adapters), ch 04-01 (Copy — why \texttt{*size\_key} is a copy, not a move).
\end{itemize}

Stop. \textbf{\textbf{Check}} — why does the gather return \texttt{Vec<PohPoint>} (owned copies of
(size, cost)) instead of \texttt{Vec<Hyperbox>} or borrows into the map? Because POH
only needs two numbers; borrowing into the map would ripple \texttt{\&mut self} rules;
and owned small copies are cheap. Return \emph{values}, not \emph{references} into
yourself — the whole point of value types.
\subsection{Step 6 — POH selection (pure transform)}
\label{sec:org7536280}

POH (potentially-optimal hyperboxes) — convex hull / pareto over the \texttt{PohPoint}
list. This is a \textbf{\textbf{pure function}}: \texttt{convex\_hull(\&[PohPoint]) -> Vec<PohPoint>}. It
touches no tree, borrows nothing lasting, allocates only its own return.

Normie-from-C++ (a method that pushes/pops a member list):
\begin{minted}[]{cpp}
void convexHull() { m_poh.push_back(...); m_poh.pop_back(); }  // mutates THIS
\end{minted}

Idiomatic (pure: takes a slice, returns an owned Vec):
\begin{minted}[]{rust}
fn convex_hull(poh: &[PohPoint]) -> Vec<PohPoint> {  // &[T], not &Vec<T>
    let mut hull: Vec<PohPoint> = Vec::new();
    for p in poh.iter() {
        while hull.len() >= 2
              && cross(hull[hull.len()-2].size, hull[hull.len()-1].size, p.size) <= 0.0 {
            hull.pop();
        }
        hull.push(*p);
    }
    hull
}
\end{minted}

Why does it feel cleaner than the C++? No side channel — input is a slice of
points, output is an owned \texttt{Vec}, nothing else changes. The absence of side
effects is not a preference; it's what lets you call it headlessly in a unit
test with a stub cost (\texttt{test-mock-traits}) with no GPU or UI. (The \texttt{\&[T]} habit
= \texttt{own-slice-over-vec} also makes it FFI-ready.)

\begin{itemize}
\item \textbf{\textbf{rust-skills}}: \texttt{anti-vec-for-slice}, \texttt{own-slice-over-vec}, \texttt{perf-collect-once},
\texttt{test-mock-traits}.
\item \textbf{\textbf{Book}}: ch 10-02 (pure funcs over traits), ch 13 (iterators), ch 04 (why an
owned \texttt{Vec<PohPoint>} return is fine — it's moved out, not copied out).
\end{itemize}
\subsection{Step 7 — trisection + the ownership dance (THE lesson)}
\label{sec:orgdfb30e2}

The step that ties the room together: take a selected box out, trisect it into
center + two shifted boxes, reinsert them. This is where Rust's ownership model
forces a structurally correct design that C++ lets you violate.

Normie-from-C++ (mutate the box in place, resize a vector):
\begin{minted}[]{cpp}
void trisect(HyperBox &b, std::vector<...> &out) { /* mutate b, push children */ }
\end{minted}

Idiomatic (consume the parent, return the children):
\begin{minted}[]{rust}
fn trisect(mut parent: Hyperbox) -> (Hyperbox, PendingTwo) {  // take = move
    let (min_idx, _) = parent.depths.iter().enumerate()
        .min_by_key(|&(_, v)| *v).expect("6-element array");
    parent.depths[min_idx] += 1;
    // center keeps parent's cost; two shifted children get cost_at_center: None
    // ... (build pos/neg shifted boxes)
    (parent, [pos, neg])
}
\end{minted}

and the caller re-inserts:
\begin{minted}[]{rust}
let (center, [pos, neg]) = Hyperbox::trisect(winner);  // winner moved OUT
// batch-eval pos/neg's new centers first, then reinsert with costs
\end{minted}

Why does Rust force ``take it out, transform, put it back''?

\begin{enumerate}
\item \textbf{\textbf{The map key is the size.}} \texttt{BTreeMap<SizeKey, ...>} — if you mutated a
map-resident box's \texttt{depths}, its \texttt{size()} changes, so its \texttt{SizeKey} would be
wrong and the map corrupted. Rust's borrow rules make this hard to even
express.
\item \textbf{\textbf{\texttt{pop\_first()} returns \emph{owned}.}} On the inner map it gives you the winner
as an owned \texttt{Hyperbox} — no lifetime to fight, free to move/transform. That's
\texttt{mem-take-replace} in the wild.
\item \textbf{\textbf{Reinsert = move=.}} \texttt{boxes.entry(new\_key).or\_default().insert(id, child)}
moves the child in. The tree \textbf{owns} it now; your Rust code no longer does.
\end{enumerate}

This is the single most important ownership lesson in the doc. Carry this:
\textbf{\textbf{the three phases map to three ownership states — borrow to gather, own to
transform, move to store — and the borrow checker keeps them honest.}}

\begin{itemize}
\item \textbf{\textbf{rust-skills}}: \texttt{own-move-large}, \texttt{mem-take-replace}, \texttt{own-copy-small},
\texttt{api-builder-must-use}.
\item \textbf{\textbf{Book}}: ch 04 (all of it); ch 15 (why \texttt{Box} fits \texttt{unique\_ptr}-shaped
ownership); ch 08 (BTreeMap entry/remove API).
\end{itemize}

Stop. \textbf{\textbf{Check}} — why can't you just mutate a box inside the map and ``put it
back in place''? Because its size (the map key) changed, and Rust won't give you
an \texttt{\&mut} into a container while you're keyed by its value (the map re-balances
on \texttt{insert=/=remove}, not in place).

That is the whole algorithm core in seven steps. What remains — cost injection,
batching, FFI, errors — is scaffolding around this loop, which is why the
Roadmap routes you back to Steps 5-7 the most.
\section{The concrete handshake}
\label{sec:org55273f9}

\subsection{One-picture flow}
\label{sec:org34f0fba}

\begin{minted}[]{text}
C++ (coordinator)                     | Rust (direct-rs)
---------------------------------------------------------------------------
RunDirectStage                        |
  backend = launch.optimizer_backend  |
  if Rust:                            |
    CppCost h = build_cpp_cost(BuildGpuCostAdapter...)   // sealed std::function
    out = direct_rs::new_direct_optimizer(range,start,budget,mode)->run(*h)
    cost_calls_=out.calls; opt_loc/val=out.best...  |  Box<DirectOptimizer>.run(&CppCost)
  else: existing DirectOptimizer path                   |       ↳ RunOutcome{best,best_cost,calls,ok}
\end{minted}
\subsection{Rust bridge (\texttt{direct-rs/src/lib.rs})}
\label{sec:orgc0add40}

\begin{minted}[]{rust}
#[cxx::bridge(namespace = "direct_rs")]
pub mod ffi {
    struct Pose { x: f64, y: f64, z: f64, xa: f64, ya: f64, za: f64 }

    enum PohMode { ConvexHull = 0, Pareto = 1, Aggressive = 2 }   // shared enum

    unsafe extern "C++" {
        include!("domain/cost.h");
        type CppCost;
        // cost hot path: 6 doubles by-value (bind to CppCost::evaluate6)
        fn evaluate6(self: &CppCost, x: f64, y: f64, z: f64,
                     xa: f64, ya: f64, za: f64) -> f64;
    }

    extern "Rust" {
        type DirectOptimizer;
        fn new_direct_optimizer(range: &Pose, start: &Pose,
                                budget: u32, mode: PohMode) -> Box<DirectOptimizer>;
        fn run(&mut self, cost: &CppCost) -> RunOutcome;     // M1: run-to-completion
        // M2 (live progress): see "Milestone options" below
    }

    // shared struct declared INSIDE the bridge so cxx generates the C++ struct
    struct RunOutcome { best: Pose, best_cost: f64, calls: u32, ok: bool }
}
\end{minted}
\subsubsection{Rust learning — what does ``shared struct'' mean, and why does cxx care?}
\label{sec:org54c9231}

The `struct Pose \{ x: f64, \ldots{} \}` inside `\#[cxx::bridge]` is a \textbf{\textbf{shared
struct}} — one definition, two languages, byte-identical layout. cxx
generates a C++ `struct Pose` with the same fields in the same order, so a
`Pose` can be passed \textbf{by value} across the boundary with zero serialization.

\begin{itemize}
\item \textbf{\textbf{Book}}: ch 05-01 (defining a struct), ch 19 (unsafe/FFI — how a `\#[repr]`
ties layout).
\item \textbf{\textbf{rust-skills}}: \texttt{api-non-exhaustive}, \texttt{name-acronym-word} (name it `Pose`,
not `POSITION6D`), \texttt{own-copy-small}. `Pose` is 6×f64 = 48 bytes — small and
`Copy`-able, so passing it by value is cheap (normie-from-Python would
over-use references/boxes worried about cost; here by-value is the fast
choice).
\end{itemize}

Normie-from-C++ (reach for a smart pointer / reference everywhere):

\begin{minted}[]{rust}
fn run(&mut self, cost: &CppCost, point: &Pose) -> f64   // borrow of a 48-byte POD
\end{minted}

Idiomatic (small value types cross by value):

\begin{minted}[]{rust}
fn evaluate6(&self, x: f64, y: f64, z: f64, xa: f64, ya: f64, za: f64) -> f64
\end{minted}

Notice the doc chose the \textbf{six-double} scalar form for the hot path, not the
`\&Pose`. That's a deliberate cxx/layout decision — a bare `Pose` by value
would also work, but six scalars make the ABI trivially C-compatible and dodge
any implicit padding question. The takeaway: \textbf{\textbf{at an FFI boundary, prefer the
simplest ABI (plain doubles / POD struct) so both compilers agree on layout
without you reasoning about padding.}} Brown Book ch 19-04 explains why layout
matters for FFI.

A shared `enum PohMode` becomes a C++ `enum class PohMode` on the other side.
Because a C++ `enum class` can legally hold an out-of-list value, every \texttt{match}
on a shared enum in Rust must include a wildcard `\_` arm — cxx cannot prove
C++ will only send a listed discriminant. This is a genuinely Rust-flavored
\textasciitilde{}normie= trap: in C++ you never match exhaustively; in Rust the compiler
\textbf{forces} the `\_` anyway.

\begin{minted}[]{rust}
let mode = match mode {
    ffi::PohMode::ConvexHull => ...,
    ffi::PohMode::Pareto     => ...,
    ffi::PohMode::Aggressive => ...,
    _ => Err(UnknownPohMode),   // cxx shared enums can't be trusted to be in-list
};
\end{minted}

(\textbf{\textbf{Book ch 06}}, exhaustive matching; this is exactly why Rust prefers exhaustiveness.)
\subsection{C++ cost-side change (\texttt{include/domain/cost.h})}
\label{sec:org5c8ead3}

Existing opaque \texttt{CppCost} wraps \texttt{std::function<double(const Point6D\&)>} and
exposes \texttt{evaluate(const Point6D\&)}. Add a 6-double bridge method so the cxx static assertion matches the Rust bridge:

\begin{minted}[]{cpp}
class CppCost {
    double evaluate(const Point6D& point) const { return fn_(point); }
    // bridge: build Point6D from 6 doubles, delegate to fn_ (no behavior change)
    double evaluate6(double x,double y,double z,
                     double xa,double ya,double za) const {
        return evaluate(Point6D(x,y,z,xa,ya,za));
    }
};
\end{minted}
\subsection{C++ call site (A/B seam) — inside \texttt{OptimizerManager::RunDirectStage} construction}
\label{sec:org6d3f999}

Scout finding: the clean kernel-level A/B point is the \textbf{DirectOptimizer
construction inside \texttt{RunDirectStage}} (\texttt{optimizer\_manager.cpp:1303}), NOT the
whole-run \texttt{OptimizerRunDriver} (driver wraps the 3-stage manager; we are
single-stage).

\begin{minted}[]{cpp}
// settings/launch: enum class OptimizerBackend { Cpp, Rust };
if (launch_.optimizer_backend == OptimizerBackend::Rust) {
    auto h = build_cpp_cost(BuildGpuCostAdapter(...)); // std::unique_ptr<CppCost>
    auto box = direct_rs::new_direct_optimizer(range,start,budget,mode);
    auto out = box->run(*h);   // *h derefs unique_ptr → CppCost& (matches &CppCost)
    cost_function_calls_      = out.calls;
    current_optimum_location_ = Point6D(out.best.x, out.best.y, /* ... */);
    current_optimum_value_    = out.best_cost;
    if (!out.ok) { emit OptimizerError("rust direct failed"); error_occurrred_ = true; }
} else {
    // existing DirectOptimizer block, byte-identical
}
\end{minted}

Rules:

\begin{itemize}
\item Do NOT put \texttt{backend} into \texttt{DirectOptimizer::Options} — ValidateOptions
fail-fast any non-default field. A settings-service field +
\texttt{OptimizerRunLaunch} is the clean home.
\item The =CppCost handle must out-live the run call (it holds the sealed
std::function capturing GPU state).
\end{itemize}
\section{Options \& landmarks}
\label{sec:orgcd1847a}

\subsection{backend selection — A/B switch mechanics}
\label{sec:org08d52fb}

\begin{enumerate}
\item \textbf{Runtime C++ virtual interface}: define an \texttt{IOptimizer} abstract seam that a
wrapper around the existing C++ \texttt{DirectOptimizer} AND the Rust backend both
implement; select at runtime (via the \texttt{OptimizerRunController} factory)
without touching the stage loop. Cleanest long-term (one objective seam the
whole app uses); more moving parts up front. cxx opaque types compose with
C++ inheritance (move-only smart-ptr-handled structs).
\item \textbf{Compile/link-time toggle} (CMake): \texttt{option(JTML\_USE\_RUST\_OPTIMIZER ...)};
conditionally \texttt{target\_link\_libraries(... direct-rs-lib)} + define
\texttt{JTML\_USE\_RUST\_OPTIMIZER}. Pick at build; zero runtime divergence; requires
rebuild to toggle.
\item \textbf{Direct single-branch swap} in \texttt{RunDirectStage} (the flow above): least
injection; both implementations coexist behind \texttt{optimizer\_backend}. Fastest
to A/B; coordination layer still knows Rust directly.
\end{enumerate}

Recommend for expedient A/B: \textbf{\textbf{\#3 first}}, then promote to a runtime interface
(\#1) if you want a settings toggle without recompiling.
\subsection{Milestones — run-to-completion vs live progress}
\label{sec:org2b12e68}

\begin{itemize}
\item \texttt{M1 (recommended first)}: \texttt{run(\&mut self, cost) -> RunOutcome} — run to
completion, no live UI progress. Shortest to a compilable A/B.
\item \texttt{M2 (live progress)}: cxx CANNOT take a C++ std::function callback into Rust
(C++→Rust fn pointers unimplemented). Reverse the wiring: an opaque C++
\texttt{ProgressSink} with methods \texttt{on\_iteration(calls)/on\_improvement(best)} that
Rust \textbf{calls out to} (Rust→C++ methods — supported) and C++ re-emits the
coordinator's \texttt{UpdateDisplay /} UpdateOptimum=. Must match C++ ordering exactly (30fps
iteration, on-improvement) or the UI regresses.
\item \texttt{M3 (per-step from C++)}: expose \texttt{step()} that C++ drives in a loop; keeps
progress control in C++ but diverges from the C++ Run structure. Not
recommended for a drop-in.
\end{itemize}
\subsubsection{Rust learning — callbacks, and why cxx flips the direction}
\label{sec:orgd845bec}

M2's \texttt{ProgressSink} may look odd if you come from C++/Python where a callback
is just ``a function you pass in''. Rust draws a sharp line between \textbf{callable
you can hand out} and \textbf{callable you can accept}, which explains the whole M2
shape.

\begin{itemize}
\item \textbf{\textbf{Book}}: ch 13-01 (closures), ch 13 (Fn-traits), ch 16 (message
passing/threads — the UI-pacing side).
\item \textbf{\textbf{rust-skills}}: \texttt{closure-fn-trait-bounds} (\texttt{FnOnce} ⊇ \texttt{FnMut} ⊇ \texttt{Fn};
request the \textbf{least} restrictive), \texttt{closure-move-capture},
\texttt{api-builder-must-use}.
\end{itemize}

The `Fn`/`FnMut`/`FnOnce` hierarchy is Rust's way of saying \textbf{how many times
and with what exclusivity} a closure may run:

\begin{itemize}
\item \texttt{FnOnce} — may be called once.
\item \texttt{FnMut} — may be called repeatedly, may mutate captured state.
\item \texttt{Fn} — may be called repeatedly, captures immutable (can be shared).
\end{itemize}

Normie-from-C++ (the host passes a \texttt{std::function} into Rust for progress):

\begin{minted}[]{cpp}
// CANNOT work in cxx: C++->Rust fn pointers aren't implemented
opt.setProgressCallback([](int iter, double best){ ... });
\end{minted}

Why cxx forbids it: function pointers in cxx flow \textbf{\textbf{Rust→C++}} only. So the
idiomatic M2 is the reverse — Rust calls a method on an opaque C++ object:

\begin{minted}[]{rust}
// Rust holds an opaque C++ ProgressSink and calls OUT (supported direction)
run_with_sink(&mut sink);          // sink.on_iteration(calls) etc.
\end{minted}

The teaching moment: in C++/Python a callback is ambient (implicit, whatever
the caller wants). In Rust a callback is a \textbf{first-class value/contract} with
well-defined effects. When the boundary forbids one direction, you invert the
ownership — the ``sink'' pattern (Rust owns a handle to the host and pokes it)
is the idiomatic Rust way to express ``the host wants progress while I run''.

M1's simpler shape (run-to-completion, return \texttt{RunOutcome}) skips callbacks
entirely — which is why it's the recommended first milestone: it lets you get
the algorithm + the cost crossing right before adding the progress-inversion
complication. This sequencing — get the pure path correct, then layer the
callback inversion — mirrors how Rust itself decouples pure logic from IO.
\subsection{Cost invocation}
\label{sec:org66a6360}

\begin{itemize}
\item \textbf{Serial per-point}: keep scalar \texttt{evaluate6(x,6doubles)} in a tight loop. Works,
but crosses the cpp boundary per point (\textasciitilde{}20k–30k total evals).
\item \textbf{Batch query} (matches C++ \texttt{BatchCostFunction}): collect new centers → one
\texttt{\&[Pose]} → \texttt{Vec<f64>} in input order. Recommended for real runs; maps onto
cxx slices.
\item \textbf{Slices}: cxx supports slices of shared structs (\texttt{Pose}) and primitives, NOT
opaque types. \texttt{\&[Pose] ↔ C++ =rust::Slice<const Pose>}. Prefer \texttt{\&[Pose]}
over a flat \texttt{Vec<f64>} for clarity.
\end{itemize}
\subsubsection{Rust learning — slices, iterators, and ``one FFI hop''}
\label{sec:org83d0c87}

The batch-vs-serial decision is a place where idiomatic Rust and performance
agree, and it teaches two habits: \textbf{\textbf{slices over references-to-Vecs}}, and
\textbf{\textbf{iterator pipelines over index loops}}.

\begin{itemize}
\item \textbf{\textbf{Book}}: ch 04-03 (slices — `\&[T]` is a fat pointer: pointer + length),
ch 13-02 (iterators), ch 13-04 (lazy iterators; collect only when needed),
ch 08 (Vec).
\item \textbf{\textbf{rust-skills}}: \texttt{own-slice-over-vec}, \texttt{perf-iter-over-index},
\texttt{perf-iter-lazy}, \texttt{perf-collect-once}, \texttt{anti-vec-for-slice},
\texttt{anti-index-over-iter}, \texttt{coll-map-choice}.
\end{itemize}

Normie-from-C++/Python (loop by index, pass \texttt{\&Vec<f64>}):

\begin{minted}[]{rust}
fn batch_cost(costs: &Vec<f64>) -> Vec<f64> {   // &Vec is a smell
    let mut out = Vec::with_capacity(costs.len());
    for i in 0..costs.len() {                    // indexing = bounds checks
        out.push(costs[i] * 2.0);
    }
    out
}
\end{minted}

Idiomatic (slice + iterator; accept `\&[f64]`, chain lazily):

\begin{minted}[]{rust}
fn batch_cost(costs: &[f64]) -> Vec<f64> {      // &[T], not &Vec<T>
    costs.iter().map(|c| c * 2.0).collect()      // lazy, no manual push
}
\end{minted}

Why this matters for the BUFFER seam specifically:

\begin{itemize}
\item \textbf{\textbf{`\&[Pose]` vs empty-receiving `Vec`}} — a slice is a \emph{view} over contiguous
memory; cxx maps `\&[T]` to C++ `rust::Slice<const T>`, which lets the C++
host pass a contiguous `std::vector<Pose>` without copying. That's the
``zero-copy batched eval'' you wanted.
\item \textbf{\textbf{One FFI hop beats many}} — each scalar `evaluate6` crossing is a boundary
transit; a batch collects \textasciitilde{}hundreds of new centers into one `\&[Pose]` and
makes \textbf{\textbf{one}} call. For 20k-30k evals this is the difference between 20k
crossings and a handful. The \texttt{BatchCostFunction} sibling in C++ exists for
precisely this.
\item \textbf{\textbf{Iterator pipelines compose with your BTreeMap}} — the POH gather you
already wrote (`self.boxes.iter()\ldots{}first\textsubscript{key}\textsubscript{value}()`) is a READ-ONLY
borrow pipeline; the trisection (`remove`/=pop\textsubscript{first}=) is the MUTATING side.
Slice-vector-boundary habits keep the borrows straight.
\end{itemize}

Reading tip: Brown Book ch 4-03's slice diagram (the fat-pointer of ptr+len)
is the exact mental model for `rust::Slice<const Pose>`.
\subsection{Cost injection — Arg-time vs constructor-time}
\label{sec:org24d74b3}

\begin{itemize}
\item \texttt{A (cost via method arg)}: \texttt{DirectOptimizer} \{ \texttt{boxes, current\_best = \} and
  =run(\&mut self, cost: \&CppCost) -> RunOutcome}. Cost flows in at call time.
\item \texttt{B (cost stored in struct)}: constructor takes the handle; Rust owns it;
\texttt{run(\&mut self)} alone. Matches the C++ constructor shape.
\end{itemize}

Both work; A keeps the optimizer pure and the FFI thin.
\subsubsection{Rust learning — the \texttt{Cost} capability as a trait (not a class)}
\label{sec:org39030e3}

Option A (cost via method arg) is the idiomatic one, and it's worth
understanding \textbf{why}: \texttt{DirectOptimizer} becomes a \textbf{\textbf{pure data + algorithm}}
object and the cost is a small \textbf{\textbf{capability}} it knows how to call, injected
per-call. In Rust that capability is best expressed as a \textbf{\textbf{trait}} — a
contract, not a base class.

\begin{itemize}
\item \textbf{\textbf{Book}}: ch 10-02 (traits), ch 17-02 (trait objects \texttt{dyn Trait}),
ch 13 (closures — the cost is eventually a closure over GPU state).
\item \textbf{\textbf{rust-skills}}: \texttt{trait-dyn-vs-generic}, \texttt{trait-object-safety},
\texttt{closure-fn-trait-bounds}, \texttt{closure-impl-fn-return}, \texttt{test-mock-traits}.
\end{itemize}

The normie-from-C++ approach (a base class \texttt{Cost} with virtual \texttt{eval()}, or a
\texttt{std::function} alias) reads as a concrete \texttt{CppCost} dependency everywhere:

\begin{minted}[]{rust}
// normie: hard-code the concrete handle into the optimizer
struct DirectOptimizer { cost: CppCost /* concrete */ }
\end{minted}

Idiomatic (a trait that \textbf{both} a test stub and the real C++ handle implement):

\begin{minted}[]{rust}
trait Cost {
    fn eval(&self, poses: &[Pose]) -> Vec<f64>;   // batch shape, needed later
}
struct CppCostHandle(/* cxx opaque */);   // impl Cost by calling evaluate6
struct AnalyticCost;                       // test stub for headless Tier-1
\end{minted}

Making \texttt{Cost} a trait vs. a concrete \texttt{CppCost} dependency is what lets you run
the SAME algorithm against a \textbf{\textbf{test stub}} and the \textbf{\textbf{real GPU cost}} — that's
\texttt{test-mock-traits}, and it's exactly why the scaffold already sketches a
\texttt{trait Cost} next to \texttt{CppCost}. Then the two shapes differ only by dispatch:

\begin{itemize}
\item \textbf{\textbf{\texttt{impl Fn} (generics)}} — zero-cost, monomorphized, great for a hot
per-call closure. Use when the cost really is a closure.
\item \textbf{\textbf{\texttt{dyn Fn=/trait object (=Box<dyn Cost>})}} — one vtable, needed when the
concrete cost is selected/stored at runtime. Your A/B already implies runtime
selection.
\end{itemize}

The A/B switch is literally this: two backends are two implementations of one
contract. Whether that contract is Rust's \texttt{Cost} trait (inside the crate) or
C++'s \texttt{IOptimizer} virtual interface (across the boundary) is the same idea in
two languages — which is what the \texttt{IOptimizer} option is pointing at.
\subsection{Outcome/error channel — Result vs status}
\label{sec:orgab86884}

\begin{itemize}
\item cxx \texttt{Result} → C++ throws \texttt{rust::Error} (std::exception). In a Qt/CUDA
pipeline, exceptions-on is uncertain → prefer \textbf{\textbf{explicit status}}:
\texttt{RunOutcome \{ ok: bool (or code), best, best\_cost, calls \}}.
\item Never let Rust panic across FFI (panic logs+aborts). Represent all error
paths in the status struct / error enum; never unwrap across the boundary.
\end{itemize}
\subsubsection{Rust learning — errors as values, not as side effects}
\label{sec:org647247e}

The whole \texttt{Result} vs status-struct decision is a Rust-native philosophy: Rust
has \textbf{\textbf{no exceptions}}. Errors are just values returned from functions. This
changes the normie-from-C++/Python default dramatically.

\begin{itemize}
\item \textbf{\textbf{Book}}: ch 09-01 (panic), ch 09-02 (\texttt{Result} + \texttt{?}),
ch 09-03 (\texttt{unwrap} vs \texttt{expect} — when each is honest).
\item \textbf{\textbf{rust-skills}}: \texttt{err-result-over-panic}, \texttt{err-thiserror-lib},
\texttt{err-anyhow-app}, \texttt{err-no-unwrap-prod}, \texttt{err-question-mark},
\texttt{err-context-chain}, \texttt{err-lowercase-msg}.
\end{itemize}

Normie-from-C++ (bool return + out-param, or throw an exception):

\begin{minted}[]{cpp}
bool Run(/* out */ double& best_cost);
\end{minted}

Normie-from-Python (rely on \texttt{raise} and \texttt{try/except}):

\begin{minted}[]{python}
def run(...):
    if bad: raise OptimizerError("bad")
\end{minted}

Idiomatic Rust (return a \texttt{Result} / explicit status; the \textbf{caller} decides):

\begin{minted}[]{rust}
type RunResult<T> = Result<T, OptimizerError>;

#[derive(thiserror::Error, Debug)]
pub enum OptimizerError {
    #[error("all-zero range {range:?}")]
    ZeroRange { range: Pose },
    #[error("non-finite cost {cost}")]
    NonFiniteCost { cost: f64 },
    #[error("poisoned executor")]
    Poisoned,
}
\end{minted}

Why the doc chose an \textbf{\textbf{explicit status struct}} (\texttt{RunOutcome\{ ok: bool, ... \}})
rather than Rust's \texttt{Result} for the FFI boundary specifically: cxx can map a
Rust \texttt{Result} to a C++ `rust::Error` exception — but that requires C++
exceptions to be \textbf{\textbf{enabled}}, which is uncertain in a Qt/CUDA toolchain, and a
Rust panic crossing an FFI boundary is \textbf{undefined behavior} (it =abort()=s the
whole process). So:

\begin{itemize}
\item \textbf{\textbf{Inside}} the Rust crate, use \texttt{Result<T, OptimizerError>} freely (=?
everywhere) — that's idiomatic and safe.
\item \textbf{\textbf{Across}} the cxx boundary, return a plain status struct, and map the inner
\texttt{Result} to it as the last step (`out.ok = result.is\textsubscript{ok}()`). This keeps the
Rust-side ergonomics and the C++-side safety.
\end{itemize}

This two-tier pattern — rich errors inside, flattened explicit status at the
FFI edge — is the idiomatic way to bridge Rust's error values to a C++ caller:

the boundary is the ONE place you convert, and you never let a panic or a
\texttt{rust::Error} exception escape across it.
\subsection{Ownership / threading}
\label{sec:org2178209}

\begin{itemize}
\item \texttt{Box<DirectOptimizer>} → C++ \texttt{rust::Box<DirectOptimizer>} (move-only,
unique\textsubscript{ptr}-like).
\item \texttt{\&self} / \texttt{\&mut self} receivers allowed; \texttt{self} by value NOT. Multi-type
extern-Rust must use explicit \texttt{self: \&T}.
\item =CppCost handle must out-live the run call (GPU captures).
\item Serial \texttt{evaluate} from one worker thread.
\end{itemize}
\subsubsection{Rust learning — ownership is the point (not a speed bump)}
\label{sec:org0c05be1}

This whole section is Brown Book ch 4 made concrete. The normie-from-C++
reaction is ``why do I need \texttt{Box<DirectOptimizer>}'', but in Rust \texttt{Box} is not a
pointer-with-extras; it's the \textbf{ownership} model made explicit.

\begin{itemize}
\item \textbf{\textbf{Book}}: ch 04-01 (what is ownership), ch 04-03 (slices),
ch 15-01 (\texttt{Box} ``for having a value on the heap''), ch 15 (smart pointers).
\item \textbf{\textbf{rust-skills}}: \texttt{own-borrow-over-clone}, \texttt{own-copy-small}, \texttt{own-move-large}.
\end{itemize}

Normie-from-C++ (\texttt{unique\_ptr} everywhere because it's the default):

\begin{minted}[]{rust}
fn make() -> Box<DirectOptimizer> { Box::new(DirectOptimizer::default()) }
\end{minted}

\ldots{}is actually fine \textbf{here} — because the C++ side owns the handle move-only, and
\texttt{Box} is Rust's \texttt{unique\_ptr}-equivalent. The lesson: \texttt{Box} isn't wrong because
it's a pointer; it's wrong to reach for it when you only need to \emph{borrow}. The
doc's \texttt{\&self} / \texttt{\&mut self} receivers are the borrow side.

Normie-from-Python mistake: passing \texttt{self} by value everywhere. cxx forbids
`self` by value for opaque types — and that's not a limitation, it's Rust
refusing to let you silently \emph{duplicate} something you can only own once. The
distinction the section hammers (Brown ch 4-01):

\begin{itemize}
\item \texttt{T} (by value): I \textbf{move} it — I alone own it now.
\item \texttt{\&T}: I \textbf{borrow} it to read; others can too.
\item \texttt{\&mut T}: I \textbf{borrow} it to mutate; exclusive.
\end{itemize}

There is no ``copy'' in that list for a non-\texttt{Copy} type. That absence is the
point. The rule ``\texttt{CppCost handle must out-live the run call}'' is likewise an
ownership statement: C++ holds the \texttt{std::function=+GPU captures; Rust borrows
it for the duration of =run}. The borrow ends when \texttt{run} returns, so the handle
stays alive — the whole lifetime idea in one sentence.
\subsection{Name / namespace}
\label{sec:org624bab5}

Use \texttt{\#[cxx::bridge(namespace = "direct\_rs")]} (or \texttt{cxx\_name=/ =rust\_name})
so \texttt{direct\_rs::DirectOptimizer ≠ ::jta::DirectOptimizer} and
\texttt{direct\_rs::Pose} stays distinct from the C++ \texttt{Point6D}.
\section{Options for the coordinate A/B seam}
\label{sec:org0944d4b}

\subsection{Seam governance the C++ side offers}
\label{sec:orgc27d162}

\begin{enumerate}
\item \texttt{OptimizerRunDriver} (factory-injected) — whole OptimizerManager lifecycle
(stages + GPU + thread); for ``alternative Manager'', not a single-stage
optimizer-core drop-in.
\item \texttt{RunDirectStage} — single-stage optimizer kernel; the least-injection point
for a Rust \texttt{DirectOptimizer} drop-in.
\item No existing switch: \texttt{direct\_options} is a \textbf{variant} slot (algorithm knobs),
not a backend selector, and fail-fast non-default fields.
\end{enumerate}
\subsection{Where \texttt{backend} state lives}
\label{sec:org80750e1}

\begin{itemize}
\item \texttt{OptimizerRunLaunch} / settings service: add \texttt{OptimizerBackend\{ Cpp, Rust \}}.
\item NOT in \texttt{Options} (fail-fast); NOT in \texttt{OptimizerRunDriver} (over-scope).
\end{itemize}
\section{Remaining open questions}
\label{sec:org4c5e7e5}

\begin{itemize}
\item batch ownership: \texttt{rust::Slice<const Pose>} over a C++ contiguous vector
(zero-copy) vs copying poses into Rust
\item \texttt{-{}-{}fno-exceptions} anywhere in Qt/CUDA toolchain? determines \texttt{Result} vs
status-struct
\item Rust \texttt{DirectOptimizer} unpin/sized; where the GPU context lives (keep single)
\item first-fixture scale (Tier-1 analytic golden) to gate Rust==C++ numeric parity
\item callback ordering for M2 (iteration-vs-im-order per eval sequence)
\end{itemize}
\section{Observations / conventions}
\label{sec:org4307221}

\begin{itemize}
\item Keep the bridge minimal; real structs/methods stay native-side; the bridge is
only value+pointer transit.
\item Re-state contract details (eval ordering, budget) here; don't rely on this
handshake alone.
\item Re-check file/line targets (\texttt{optimizer\_manager.cpp:1303}, \texttt{cost.h}, \texttt{lib.rs})
before coding — they drift.
\item Verify verbatim generated signatures (rust::Box, ::Slice) with a tiny smoke
crate first; don't trust this doc's shapes.
\end{itemize}
\section{Review policy}
\label{sec:org2e0540b}

This is a \textbf{reference handshake}, not an implementation plan. Decisions to
double-check before coding:

\begin{enumerate}
\item \texttt{evaluate6} 6-double hot path vs shared-struct batch seam — which first?
\item A/B seam = RunDirectStage construction vs \texttt{IOptimizer} virtual adapter.
\item status-struct return vs cxx \texttt{Result} (exception).
\item M1-only vs live progress day-one for the first A/B test.
\end{enumerate}

(Note: exact generated-header shapes for \texttt{rust::Box<T>} and \texttt{rust::Slice<const T>} were
not verified verbatim from a build artifact; cxx latest \textasciitilde{}1.0.194. Build a tiny smoke)
\section{Appendix A — deep dive: borrowing and the DIRECT loop}
\label{sec:org3a7c0b6}

If you only read one appendix, make it this one: it's ownership/borrowing made
concrete against the exact shapes this project uses. The Brown Book ch 4 is the
source; here we map it onto \texttt{boxes: BTreeMap<SizeKey, BTreeMap<CostKey,
Hyperbox>{}>{}} and the run loop.
\subsection{The four borrow flavors and when the loop uses each}
\label{sec:orgd12017e}

Re-read this table whenever a borrow error confuses you. Every one of these is
deliberately exercised by the DIRECT loop.

\begin{center}
\begin{tabular}{lll}
Rust expression & borrow kind & in this project\\
\hline
\texttt{\&self} / \texttt{\&T} & shared (read-only) & \texttt{gather\_poh} over the whole tree\\
\texttt{\&mut self} / \texttt{\&mut T} & exclusive (write) & \texttt{trisect\_and\_reinsert} (it removes + inserts)\\
\texttt{T} by value & move (ownership) & \texttt{pop\_first()} giving you the owned winner\\
\texttt{\&[T]} & slice (a window) & \texttt{convex\_hull(\&[PohPoint])} — a borrowed view\\
\end{tabular}
\end{center}

The rule that governs almost every error here: \textbf{\textbf{you can have many \texttt{\&T} or one
\texttt{\&mut T}, never both for the same data at the same time.}} (Brown ch 4-02.)
\subsection{The borrow dance, zoomed to one iteration}
\label{sec:org7e223e8}

Walk the loop and name the borrow at each breakpoint (this is Milestone 9's
self-check, but written out):

\begin{minted}[]{rust}
// 1. GATHER — shared borrow: many readers, no writer
let mut candidate_points = self.gather_poh();     // &self
let selected = make_hull(&candidate_points);       // &[PohPoint]

// 2. TRISECT — here we switch from &self to &mut self
//    `self` was &self during gather; now we need &mut self to mutate the map.
for point in &selected {
    // 3. TAKE — move ownership OUT of the map (needs &mut)
    let winner: Hyperbox = {
        let row = self.boxes.get_mut(&OrderedFloat(point.size)); // &mut row
        let (_key, hb) = row.expect("selected row exists").pop_first()
            .expect("non-empty row");
        hb                                          // moves the Hyperbox out
    }; // row borrow ends here

    // 4. TRANSFORM — we OWN `winner` now; no borrow at all
    let (center, [pos, neg]) = Hyperbox::trisect(winner);

    // 5. REINSERT — move the children BACK in (&mut)
    self.insert_box(center);   // insert as known-cost
    // pos/neg: batch-eval their centers, then insert with scores
}
\end{minted}

The borrow transitions are the \textbf{spine} of the method:

\begin{itemize}
\item \texttt{\&self→\&mut self} happens by \textbf{method signature}, not by hand — a method
\end{itemize}
either takes \texttt{\&self} or \texttt{\&mut self}, so the compiler forces gather and
trisect to be separate methods with different receivers. That's why you can't
have one giant \texttt{\&mut} method that gathers and mutates at once and expect the
borrow rules to stay sane.
\begin{itemize}
\item The \texttt{row} borrow in step 3 is \textbf{\textbf{scoped to the block}}: the \texttt{\&mut row} dies
\end{itemize}
at the closing brace, so \texttt{hb} (owned) can be returned out of it. This block-
scoping is exactly the \texttt{mem-take-replace} / last-use pattern; the compiler ends
the borrow at the winner's last use.
\subsection{The three errors you will actually hit (the canon)}
\label{sec:org835c550}

These are the borrow-error codes you'll encounter in this project, mapped to
the \emph{cause} and the \emph{fix}:

\begin{center}
\begin{tabular}{llll}
Error & What it means & How it bites here & The fix\\
\hline
\texttt{E0382} & use of a moved value: you moved it away & reusing a \texttt{Hyperbox} after & don't reuse after move; or\\
 & then used it again & \texttt{pop\_first()} moved it out & reorder so last use comes first\\
\texttt{E0502} & cannot borrow as mutable because it's also & \texttt{\&mut self} while a \texttt{\&self} & scope the shared borrow so it ends\\
 & borrowed as immutable (two borrows, wrong way) & borrow of the map is still alive & before the \texttt{\&mut} begins\\
\texttt{E0499} & two mutable borrows at once & two live \texttt{\&mut} into the tree & never hold two \texttt{\&mut} into the\\
 &  & (e.g. nested \texttt{get\_mut} loops) & same map; restructure\\
\texttt{E0505} & cannot move a value out of a borrow & \texttt{\&self} row then trying to & pop/remove \emph{out} via \texttt{\&mut} or\\
 &  & return an owned \texttt{Hyperbox} & clone only the fields you need\\
\end{tabular}
\end{center}

\begin{itemize}
\item \textbf{\textbf{Book}}: the Brown Book has a whole chapter (ch 4) that animates these,
with an interactive ``editor'' where you drag borrows around.
\item \textbf{\textbf{rust-skills}}: \texttt{own-borrow-over-clone}, \texttt{own-move-large},
\texttt{mem-take-replace}, \texttt{anti-index-over-iter}.
\end{itemize}

Stop. \textbf{\textbf{Check}} — you hit ``cannot move out of \texttt{*cost} which is behind a \texttt{\&
reference" earlier. Which error code is that (E0505), and what's the conceptual
fix — borrow the /fields/ you need, or move the whole value out via =\&mut}? It
maps to ''project the two numbers (Copy fields) rather than pulling the whole
box."
\section{Appendix B — deep dive: ownership, lifetimes, and Send/Sync}
\label{sec:orgcd07793}

This appendix is for the moments when value/move/borrow isn't enough — when you
need \texttt{Box}, when a lifetime shows up in a signature, and when the compiler
starts asking about \texttt{Send} and \texttt{Sync} because of the worker thread + GPU cost.
\subsection{Lifetimes: what they really are (and aren't)}
\label{sec:orgd25fa93}

A lifetime is not a value you pass around; it's the compiler's label for ``how
long a borrow is valid''. You mostly write \texttt{\&T} and let Rust elide lifetimes; an
explicit lifetime (\texttt{<'a>}) appears when a return type borrows from an argument.

Normie-from-C++ confusion: ``\texttt{\&Point6D} is a const reference, why do I need a
lifetime?'' The answer: the returned value's lifetime is tied to the borrow's
origin. In this project you rarely write explicit lifetimes because you \emph{own}
the values. The one place they appear is cxx opaque types and \texttt{\&CppCost}
borrows — which cxx manages for you.

\begin{itemize}
\item \textbf{\textbf{Book}}: ch 10-03 (lifetimes), Brown ch 4-02 (the borrow table).
\item \textbf{\textbf{rust-skills}}: \texttt{own-lifetime-elision}, \texttt{type-phantom-marker}.
\end{itemize}
\subsection{\texttt{Box}: only when you need a stable address or heap}
\label{sec:orgd5f35c1}

Recall \texttt{rust::Box<DirectOptimizer>} on the C++ side. Box is one specific tool,
not a default. You reach for \texttt{Box<T>} when:

\begin{itemize}
\item T is recursive (a tree node) — not here.
\item T is huge and you want to move a pointer instead — your \texttt{Hyperbox} is tiny.
\item you need heap because of a trait object \texttt{Box<dyn Cost>} — only if you pick
runtime dispatch for the cost (Milestone 6).
\item FFI wants a stable pointer (cxx \texttt{Box<DirectOptimizer>}) — THIS is your case.
\end{itemize}

Normie mistake: \texttt{Box} every struct ``just in case''. Stop — value \texttt{Hyperbox=`s by
default; box only the FFI boundary object. =own-move-large}, \texttt{mem-box-large-
variant}.
\subsection{Send and Sync: why the GPU cost blocks one thread}
\label{sec:org3838b35}

This is where \texttt{Send} and \texttt{Sync} finally matter. The worker thread (on the C++
side) calls \texttt{run} synchronously; the GPU cost (:CppCost=) mutates \texttt{GPUModel}
pose state on each \texttt{evaluate}. So:

\begin{itemize}
\item \texttt{Send} means ``value can be transferred to another thread.'' \texttt{CppCost} wraps
a \texttt{std::function} + GPU captures — Rust can't prove it's thread-safe, so it
won't be auto-\texttt{Send}.
\item \texttt{Sync} means ``\&T can be shared across threads.'' Same story.
\end{itemize}

Rather than fight that, the design keeps everything on ONE worker thread and
calls \texttt{run} to completion (M1). That sidesteps \texttt{Send} / \texttt{Sync} entirely — you never
move \texttt{CppCost} between threads because the C++ side owns the thread and the
handle. This is why the doc's threading rule is ``serial, one thread'' and why
=M1 (run-to-completion) is recommended before M2 (which needs care about which
thread calls into the sink).

\begin{itemize}
\item \textbf{\textbf{Book}}: ch 16 (concurrency — Send/Sync are the runtime of fear \texttt{free}
threading; \texttt{Rc} vs \texttt{Arc}).
\item \textbf{\textbf{rust-skills}}: \texttt{own-arc-shared} (when you DO need shared ownership across
threads, upgrade \texttt{Rc=→=Arc}), \texttt{conc-thread-local}, \texttt{unsafe-send-sync-manual}.
\end{itemize}
\subsection{\texttt{Rc=/=Arc=/=RefCell}: the ``when you DON'T need them'' warning}
\label{sec:org27e1f61}

If you find yourself reaching for \texttt{Rc}, \texttt{Arc}, or \texttt{RefCell} in DIRECT, pause
hard. They are for \emph{shared} or \emph{interior} ownership, and the whole design avoids
both:

\begin{itemize}
\item \texttt{Rc/Arc} for shared ownership — you never share a \texttt{Hyperbox}; the map owns it
and you move it in/out. You don't need shared pointers at all.
\item \texttt{RefCell} for interior mutability — you mutate through \texttt{\&mut} methods, not
through interior mutation. If you're grabbing \texttt{RefCell}, you've probably
mis-structured a method's receiver (should be \texttt{\&mut self}, not \texttt{\&self} + a
hidden \texttt{RefCell}).
\end{itemize}

Intro Section n/a; see Interior Mutability appendix.

Stop. \textbf{\textbf{Check}} — why is \texttt{CppCost} not auto-\texttt{Send/Sync}, and how does keeping
the run on one thread dodge that (instead of adding \texttt{unsafe impl Send})? Because
Rust can't verify the captured GPU state is thread-safe; single-threading makes
the question moot without any unsafe.
\section{Appendix C — interior mutability: when you do NOT need it}
\label{sec:org7b1ed64}

Interior mutability (\texttt{Cell}, \texttt{RefCell}, \texttt{Mutex}, \texttt{RwLock}, \texttt{Atomics}) lets you
mutate through a shared reference. It is a \textbf{workaround for a borrow it gives
up}, and DIRECT's design exists largely to not need it. This appendix helps you
recognize the moment you're reaching for it and whether that's a real need or a
sign you've mis-structured.
\subsection{The honest use case}
\label{sec:org103548a}

You need interior mutability when: a value must be mutably changed, but you
only have \texttt{\&self=/=\&} access, and you can't re-structure to \texttt{\&mut}. Classic
examples: caches, counters behind an API that only hands out \texttt{\&self}, or raw
pointer mutation like GPU buffers owned by C++ and visible to Rust as \texttt{\&T}.

In DIRECT, the places you might 'tempt' reach for it:

\begin{itemize}
\item a \texttt{cost\_function\_calls\_} counter you bump in \texttt{\&self} methods;
\item a cache of computed \texttt{size()} values;
\item the \texttt{CppCost} handle that mutates GPU state but comes in as \texttt{\&CppCost}.
\end{itemize}
\subsection{Why DIRECT mostly dodges it}
\label{sec:org307aab7}

\begin{enumerate}
\item \textbf{\textbf{The counter.}} In the C++ it's a member bumped by all methods. The Rust
idiomatic move is to make \texttt{\&mut}-taking methods the ONLY mutators, and keep
the counter next to \texttt{boxes} behind \texttt{\&mut self}. If a counter needs to be
readable from \texttt{\&self} AND written from elsewhere, \texttt{Cell<u32>} is legitimate
— but only because you've proven \texttt{\&self} must observe it while \texttt{\&mut}
writes elsewhere. Often you can reorder so a single \texttt{\&mut self} method owns
the bump and the read, and \texttt{Cell} never appears.
\item \textbf{\textbf{The cache.}} \texttt{perf-profile-first}: don't cache until the profiler says so.
And even then, a \texttt{RefCell} cache turns every read into a potential panic
(.borrow() twice = panic) — you trade one class of bug for another.
\item \textbf{\textbf{The GPU \texttt{CppCost}.}} This is the real one. On the C++ side, \texttt{evaluate()}
mutates \texttt{GPUModel} pose via a mutable closure, but to Rust it arrives as
\texttt{\&CppCost}. If Rust needed to \emph{store} that handle and call it from multiple
places, you might reach for \texttt{RefCell<\&CppCost>}. But the design passes =
\&CppCost= /in to run/=/batch= and finishes the call synchronously, so the
C++ side alone owns mutation — no interior mutability needed in Rust.
\end{enumerate}
\subsection{The ``is it a real need?'' checklist **}
\label{sec:org9d9a47f}

Ask, in order, before introducing \texttt{RefCell} / \texttt{Cell}:

\begin{enumerate}
\item Can I change a \texttt{\&self} method to \texttt{\&mut self}? If a mutator touches state,
it should \emph{take} \texttt{\&mut} — that's the honest signal.
\item Can I return data instead of mutating in place? (gather returns \texttt{Vec<PohPoint>};
it doesn't scribble a member.)
\item Is this genuinely shared mutable state visible across \& reads and \&mut
writes? Only \emph{then} is \texttt{Cell=/=RefCell=/=Mutex} justified, and pick \texttt{Mutex}
if it crosses threads.
\end{enumerate}

If the answer to \#1 or \#2 is yes, you don't need interior mutability — you need
a better method signature. Reaching for \texttt{RefCell} early is the ``masturbatory
types'' spike this project wants to avoid (see \texttt{anti-over-abstraction} and the
Options-as-data block).

\begin{itemize}
\item \textbf{\textbf{Book}}: ch 15-05 (\texttt{RefCell} ``interior mutability''), ch 16 (threads =
\texttt{Mutex} vs \texttt{RwLock}), rust-skills \texttt{own-refcell-interior} =
\end{itemize}
\texttt{own-mutex-interior}.
\begin{itemize}
\item \textbf{\textbf{rust-skills}}: \texttt{own-refcell-interior}, \texttt{own-mutex-interior},
\texttt{own-rc-single-thread}, \texttt{anti-over-abstraction}.
\end{itemize}

Stop. \textbf{\textbf{Check}} — you're about to add \texttt{RefCell<u32>} for a call counter. Run
the checklist: can the bumping method take \texttt{\&mut self}? If yes, drop the
\texttt{RefCell}. Only keep it when \texttt{\&self} readers MUST see the current count while
another object owns mutation — which, in DIRECT, the \texttt{\&mut self}-run loop
handles without interior mutability.
\section{Appendix D — \texttt{Option} and \texttt{Result} in practice (the Python jump)}
\label{sec:org82575e4}

Coming from Python, this is the appendix that will save you the most time.
Python has \texttt{None} and exceptions and they're \textbf{ambient}: any value might be
\texttt{None}, any call might raise. Rust makes both into ordinary types you must open
before use. \texttt{Option<f64>} is not ``an f64 that might be missing'' — it is a
different type from \texttt{f64}, and the compiler will not let you add 1.0 to it.

This is exactly the \texttt{cost\_at\_center: Option<f64>} field from Step 3, so every
example below is real.

\begin{itemize}
\item \textbf{\textbf{Book}}: ch 06-02 (\texttt{match} + \texttt{Option}), ch 06-03 (\texttt{if let}),
ch 09-02 (\texttt{Result} + \texttt{?}), ch 18 (patterns in depth).
\item \textbf{\textbf{rust-skills}}: \texttt{type-option-nullable}, \texttt{err-question-mark},
\texttt{err-no-unwrap-prod}, \texttt{pat-exhaustive-enum}.
\end{itemize}
\subsection{The four ways to open an \texttt{Option} (in order of preference)}
\label{sec:orgab00a42}

\begin{minted}[]{rust}
// 1. match — when both arms do real work
match hb.cost_at_center {
    Some(c) => record(c),
    None    => queue_for_eval(hb),
}

// 2. if let — when only one arm does work
if let Some(c) = hb.cost_at_center { record(c); }

// 3. let-else — when None means "bail out of this function" (the flattener)
let Some(c) = hb.cost_at_center else { return Err(OptimizerError::Unscored) };
// `c` is a plain f64 from here down — no nesting, no indentation creep

// 4. combinators — when you're transforming, not branching
let doubled: Option<f64> = hb.cost_at_center.map(|c| c * 2.0);
let shown:   f64         = hb.cost_at_center.unwrap_or(f64::INFINITY);
\end{minted}

\texttt{let-else} (option 3) is the one C++/Python developers never discover on their
own and then use constantly. It is the idiomatic early-return: the happy path
stays un-indented, the failure path is one line. Reach for it whenever your
\texttt{match} has a \texttt{None} arm that just leaves.
\subsection{The combinator cheat-sheet}
\label{sec:org730c7f9}

These are the \texttt{Option=/=Result} methods worth memorizing. Every one of them
replaces an \texttt{if} you'd have written in C++.

\begin{center}
\begin{tabular}{lll}
Call & Means & Use here\\
\hline
\texttt{.map(f)} & transform the inside, keep the wrapper & scale a cost, keep ``unscored''-ness\\
\texttt{.and\_then(f)} & \texttt{map} where \texttt{f} itself returns Option/Result & chained lookups (row → box)\\
\texttt{.unwrap\_or(v)} / \texttt{\_else(f)} & supply a default & treat unscored as \texttt{INFINITY} for sorts\\
\texttt{.ok\_or(e)} / \texttt{\_else(f)} & \texttt{Option} → \texttt{Result} (attach a reason) & ``selected row missing'' → error\\
\texttt{.ok()} / \texttt{.err()} & \texttt{Result} → \texttt{Option} (throw the reason away) & the FFI-edge flatten\\
\texttt{.filter(p)} & keep \texttt{Some} only if predicate holds & ``scored AND finite''\\
\texttt{.take()} & move the value out, leave \texttt{None} behind & claim a pending cost exactly once\\
\texttt{.as\_ref()} / \texttt{.as\_mut()} & \texttt{\&Option<T>} → \texttt{Option<\&T>} & peek without moving out of \texttt{\&self}\\
\texttt{.is\_some\_and(p)} & test the inside without unwrapping & guards in \texttt{if}\\
\texttt{.zip(other)} & pair two \texttt{Some} into one \texttt{Some((a,b))} & need both center costs or neither\\
\end{tabular}
\end{center}

\texttt{.take()} deserves a flag: it is \texttt{mem-take-replace} in miniature and it's how
you move a non-\texttt{Copy} value out of a struct field you only have \texttt{\&mut} to.
When the borrow checker says ``cannot move out of \texttt{self.field}'', \texttt{.take()} is
usually the answer.
\subsection{\texttt{?} — the operator that makes \texttt{Result} bearable}
\label{sec:org287c578}

Normie (Python habit: let it raise, or C++ habit: check every return):

\begin{minted}[]{rust}
let row = match self.boxes.get_mut(&key) {
    Some(r) => r,
    None => return Err(OptimizerError::MissingRow { key }),
};
let (_k, hb) = match row.pop_first() {
    Some(p) => p,
    None => return Err(OptimizerError::EmptyRow { key }),
};
\end{minted}

Idiomatic (\texttt{?} = ``unwrap or return the error to my caller''):

\begin{minted}[]{rust}
let row = self.boxes.get_mut(&key).ok_or(OptimizerError::MissingRow { key })?;
let (_k, hb) = row.pop_first().ok_or(OptimizerError::EmptyRow { key })?;
\end{minted}

Six lines to one, with the \textbf{same} control flow and no hidden unwinding. This is
what ``errors are values'' buys you: \texttt{?} is sugar for an early return, nothing
more. It works on \texttt{Option} too (returns \texttt{None}), which is why small helper
functions often return \texttt{Option<T>} just to get \texttt{?}.
\subsection{\texttt{unwrap} vs \texttt{expect} vs neither}
\label{sec:org6bac710}

\begin{itemize}
\item \texttt{.unwrap()} — never in this crate's non-test code. It's a panic, and
Appendix B tells you a panic across FFI aborts the process.
\item \texttt{.expect("6-element array")} — acceptable \textbf{\textbf{only}} when the invariant is
structurally guaranteed, and the string says \textbf{why}. The \texttt{min\_by\_key} call in
Step 7 is the legitimate case: a \texttt{[u32; 6]} is never empty.
\item Anything else — \texttt{?} or an explicit \texttt{match}.
\end{itemize}

Stop. \textbf{\textbf{Check}} — you have \texttt{Option<f64>} and need an \texttt{f64} to compare. Name
three different correct moves and when each is right. (\texttt{unwrap\_or(INFINITY)} if
unscored should sort last; \texttt{let-else} if unscored is a bug worth returning on;
\texttt{is\_some\_and(|c| c < best)} if you're only testing.)
\section{Appendix E — the run loop assembled (contract-faithful)}
\label{sec:org53fd74e}

Milestone 9 says ``wire it together'' and the doc never shows the wiring. Here it
is, and the reason it's worth spelling out is that \textbf{\textbf{the eval-ordering contract
from ``Contract being filled'' lives in exactly one function}} — get it wrong and
the Tier-1 golden diverges in a way that's miserable to bisect.

Re-read the contract line first:

> eval ordering: calls++ → non-finite → optimum update → storage → improvement callback
\subsection{The struct}
\label{sec:orgf141497}

\begin{minted}[]{rust}
pub struct DirectOptimizer {
    tree: BoxTree,                 // the two-level BTreeMap (Step 4)
    budget: u32,                   // cumulative cap: trunk 20k → branch 25k/30k → leaf 35k
    call_offset: u32,              // SetCallOffset — cumulative across stages
    calls: u32,
    non_finite: u32,
    best: Pose,
    best_cost: f64,
    stop: Arc<AtomicBool>,         // cooperative Stop(), set from the C++ thread
}
\end{minted}

That \texttt{Arc<AtomicBool>} is the \textbf{one} honest interior-mutability case in the
whole design, and it's worth pausing on — Appendix C told you to be suspicious
of \texttt{RefCell}, and this is the exception that proves the rule. \texttt{Stop()} is
called from a different thread than \texttt{run()}, so there is no way to restructure
it into \texttt{\&mut self}: the checklist's question 3 (``genuinely shared mutable
state across threads?'') answers \textbf{yes}. \texttt{AtomicBool} not \texttt{Mutex<bool>} because
it's one bit and never needs to be read-and-written atomically together with
anything else. Everything else in the struct stays behind \texttt{\&mut self}.
\subsection{The eval helper — where the contract is enforced}
\label{sec:org9e959f9}

\begin{minted}[]{rust}
/// Evaluate a batch and fold the results into optimizer state.
/// Returns costs in INPUT ORDER (the contract the C++ BatchCostFunction keeps).
fn eval_batch(&mut self, cost: &dyn Cost, poses: &[Pose]) -> Vec<f64> {
    let costs = cost.eval(poses);            // ONE FFI hop for the whole batch
    debug_assert_eq!(costs.len(), poses.len(), "batch cost must preserve length");

    for (pose, &c) in poses.iter().zip(costs.iter()) {
        self.calls += 1;                     // 1. calls++  (counts non-finite too)
        if !c.is_finite() {
            self.non_finite += 1;            // 2. non-finite
            continue;                        //    a non-finite eval never becomes best
        }
        if c < self.best_cost {              // 3. optimum update
            self.best = *pose;
            self.best_cost = c;
            self.improved_this_iter = true;  // 5. deferred — see below
        }
    }
    costs                                    // 4. storage is the caller's job (reinsert)
}
\end{minted}

The subtle bit: the contract puts \textbf{\textbf{storage before the improvement callback}},
but the improvement is \textbf{detected} here during the eval. So you don't fire the
sink here — you set a flag and fire it after reinsertion. Firing early is the
classic port bug: the UI shows a new optimum for a box that isn't in the tree
yet, and if anything downstream reads the tree on that signal, it races.
\subsection{The loop}
\label{sec:orgcb7b416}

\begin{minted}[]{rust}
pub fn run(&mut self, cost: &dyn Cost, sink: Option<&mut dyn ProgressSink>) -> RunResult<()> {
    if self.tree.is_empty() { return Err(OptimizerError::ZeroRange { range: self.range }); }

    while (self.calls + self.call_offset) < self.budget
          && !self.stop.load(Ordering::Relaxed)
    {
        self.improved_this_iter = false;

        // GATHER — &self borrow, no mutation (Step 5)
        let points   = self.tree.gather_poh();
        let selected = convex_hull(&points);          // pure (Step 6)
        if selected.is_empty() { break; }             // nothing left to split

        // TRISECT — &mut self; take out, transform, put back (Step 7)
        let mut pending: Vec<(Pose, PendingSlot)> = Vec::new();
        for point in &selected {
            let winner = self.tree.take_min_of_row(point.size)?;   // owned
            let (center, [pos, neg]) = Hyperbox::trisect(winner);
            self.tree.insert_scored(center);                       // cost inherited
            pending.push((pos.center, PendingSlot::Pos(pos)));
            pending.push((neg.center, PendingSlot::Neg(neg)));
        }

        // ONE FFI hop for every new center this iteration
        let poses: Vec<Pose> = pending.iter().map(|(p, _)| *p).collect();
        let costs = self.eval_batch(cost, &poses);

        // STORAGE — then and only then, the callbacks
        for ((_, slot), c) in pending.into_iter().zip(costs) {
            self.tree.insert_with_cost(slot.into_box(), c);
        }
        if let Some(s) = sink.as_deref_mut() {
            s.on_iteration(self.calls + self.call_offset);   // ~30fps pacing is C++'s job
            if self.improved_this_iter { s.on_improvement(self.best, self.best_cost); }
        }
    }
    Ok(())
}
\end{minted}

And the FFI edge — the one place \texttt{Result} becomes status (M8):

\begin{minted}[]{rust}
pub fn run_ffi(&mut self, cost: &CppCost) -> ffi::RunOutcome {
    let handle = CppCostHandle(cost);
    let result = self.run(&handle, None);
    ffi::RunOutcome {
        best:      self.best,
        best_cost: self.best_cost,
        calls:     self.calls,
        ok:        result.is_ok(),      // rich error → one bool, deliberately
    }
}
\end{minted}
\subsection{Three things this skeleton exposes that the prose didn't}
\label{sec:orgf09a45a}

\begin{enumerate}
\item \textbf{\textbf{Batch granularity is per-iteration, not per-box.}} All the new centers
from one POH round go in one \texttt{\&[Pose]}. That's the ``one FFI hop'' claim made
concrete: \textasciitilde{}2×|hull| poses per crossing instead of one.
\item \textbf{\textbf{The budget can overshoot.}} \texttt{eval\_batch} evaluates the whole batch, so the
final iteration may push \texttt{calls} past \texttt{budget}. The C++ checks the guard
\textbf{between} iterations too — so this matches — but confirm it against
\texttt{direct\_optimizer.h}, because if C++ truncates mid-batch and Rust doesn't,
the golden diverges on call count. \textbf{\textbf{This belongs in ``Remaining open
questions''.}}
\item \textbf{\textbf{\texttt{sink} is \texttt{Option<\&mut dyn ProgressSink>}.}} M1 passes \texttt{None} and the
whole progress path costs nothing; M2 passes \texttt{Some}. One signature, both
milestones — this is \texttt{trait-dyn-vs-generic} paying off.
\end{enumerate}

Stop. \textbf{\textbf{Check}} — trace the calls counter for one iteration with a 4-box hull.
How many FFI crossings, how many \texttt{calls} increments, and in what order does a
new best reach the UI?
\section{Appendix F — testing the port: unit, property, parity}
\label{sec:org30a290e}

A port has one job: \textbf{be the same}. That makes this the appendix with the
highest ratio of value to Rust-novelty, because Rust's test story is unusually
good and you can get parity confidence before any C++ is linked.

\begin{itemize}
\item \textbf{\textbf{Book}}: ch 11 (all of it — it's short), ch 11-03 (integration tests).
\item \textbf{\textbf{rust-skills}}: \texttt{test-mock-traits}, \texttt{err-result-over-panic},
\texttt{perf-profile-first}.
\end{itemize}
\subsection{Tier 0 — unit tests live next to the code}
\label{sec:orgd955d62}

\begin{minted}[]{rust}
// bottom of src/poh.rs
#[cfg(test)]                     // compiled ONLY under `cargo test`
mod tests {
    use super::*;

    #[test]
    fn hull_of_a_known_set() {
        let pts = vec![p(1.0, 5.0), p(2.0, 3.0), p(3.0, 4.0), p(4.0, 1.0)];
        let hull = convex_hull(&pts);
        assert_eq!(hull.len(), 3);
        assert!(hull.windows(2).all(|w| w[0].size < w[1].size));  // monotone in size
    }

    #[test]
    fn trisect_preserves_depth_budget() {
        let parent = Hyperbox { cost_at_center: Some(1.0), center: Pose::ZERO, depths: [0; 6] };
        let (center, [pos, neg]) = Hyperbox::trisect(parent);
        assert_eq!(center.depths.iter().sum::<u32>(), 1);
        assert!(pos.size() < 1.0 && neg.size() < 1.0);
        assert!(pos.cost_at_center.is_none());     // not scored yet — Step 3's Option
    }
}
\end{minted}

The \texttt{\#[cfg(test)] mod tests} idiom is worth internalizing: tests sit in the
\textbf{same file} as the code and can see private items. There is no separate test
tree to keep in sync, and no header to expose internals through. This is one of
the places where Rust is simply nicer than C++ and you should exploit it —
every Step in this doc should end with a test module.
\subsection{Tier 1 — the analytic golden, headless}
\label{sec:org59b7e96}

This is why \texttt{Cost} is a trait (M6). With a stub cost you can run the \textbf{entire}
algorithm with no GPU, no Qt, no C++:

\begin{minted}[]{rust}
struct Sphere;                                  // f(x) = sum(x_i^2), min at origin
impl Cost for Sphere {
    fn eval(&self, poses: &[Pose]) -> Vec<f64> {
        poses.iter()
             .map(|p| p.x*p.x + p.y*p.y + p.z*p.z + p.xa*p.xa + p.ya*p.ya + p.za*p.za)
             .collect()
    }
}

#[test]
fn finds_the_origin_within_budget() {
    let mut opt = DirectOptimizer::new(range(), start(), 20_000, PohMode::ConvexHull);
    opt.run(&Sphere, None).unwrap();
    assert_relative_eq!(opt.best_cost, 0.0, epsilon = 1e-6);
    assert!(opt.calls <= 20_000);
}
\end{minted}
\subsection{Tier 2 — property tests for the invariants prose can't hold}
\label{sec:org0ad58c2}

\texttt{proptest} generates hundreds of inputs and shrinks failures to a minimal case.
Use it for the claims this doc makes in English:

\begin{minted}[]{rust}
proptest! {
    // the hull is always a subset of its input
    #[test]
    fn hull_is_subset(pts in prop::collection::vec(any_poh_point(), 1..50)) {
        for h in convex_hull(&pts) { prop_assert!(pts.contains(&h)); }
    }

    // trisection never loses volume: children's extents tile the parent's
    #[test]
    fn trisect_conserves(depths in any::<[u32; 6]>()) {
        let parent = Hyperbox { cost_at_center: Some(0.0), center: Pose::ZERO, depths };
        let (c, [p, n]) = Hyperbox::trisect(parent);
        prop_assert_eq!(c.depths.iter().sum::<u32>(),
                        depths.iter().sum::<u32>() + 1);
        prop_assert!(p.size() <= parent_size && n.size() <= parent_size);
    }
}
\end{minted}

Property testing is not a Rust-specific idea, but Rust's \texttt{proptest} is where
most people meet it, and ``the same optimum for the same cost sequence'' is
precisely a property, not an example.
\subsection{Tier 3 — parity against the C++ (the actual gate)}
\label{sec:org0f5e47e}

The contract says the judge is ``the SAME optimum for the same cost sequence''.
So make the cost sequence the artifact:

\begin{enumerate}
\item Instrument the C++ \texttt{DirectOptimizer} to dump, per eval: index, the 6 doubles
in, the cost out. Run one fixture. That file is your golden.
\item In Rust, implement \texttt{Cost} as a \textbf{replay} over that file: it asserts the
incoming pose matches the recorded pose and returns the recorded cost.
\item Any divergence fails at the \textbf{first differing eval}, with an index — instead
of you comparing two final optima and having no idea where they split.
\end{enumerate}

\begin{minted}[]{rust}
struct ReplayCost { records: Vec<(Pose, f64)>, cursor: Cell<usize> }
// eval(): for each pose, assert_relative_eq! against records[cursor], return cost
\end{minted}

This is the single highest-leverage test in the port, and note what it needs:
\texttt{Cost} being a trait, and the batch/serial distinction being a detail the trait
hides. Two design decisions from M6 exist mostly to make this test possible.

(\texttt{Cell<usize>} for the cursor is legitimate here — Appendix C's checklist,
question 3, with \texttt{\&self} forced on you by the trait signature. A test is also
the one place a \texttt{RefCell} panic is acceptable: it fails the test, which is the
point.)

Stop. \textbf{\textbf{Check}} — which tier catches ``Rust picks a different tie-break when two
boxes have equal cost''? (Tier 3, immediately, at the first tie — and Tier 2 if
you write the property. Tier 1 might never notice, because the final optimum
can match by luck.)
\section{Appendix G — C++ → Rust and Python → Rust translation tables}
\label{sec:orgb8611bf}

A lookup for the ``what's the Rust word for this'' moment. Use it to stop
searching, not to stop thinking — the third column is where the difference
lives.
\subsection{From C++}
\label{sec:org1cfde3b}

\begin{center}
\begin{tabular}{lll}
C++ & Rust & The difference that matters\\
\hline
\texttt{std::vector<T>} & \texttt{Vec<T>} & \texttt{Vec} is move-only unless \texttt{T: Clone}; no copy ctor\\
\texttt{const std::vector<T>\&} arg & \texttt{\&[T]} & slice works for arrays/Vec/boxed alike\\
\texttt{std::array<T, N>} & \texttt{[T; N]} & size is part of the type, same as C++\\
\texttt{std::unique\_ptr<T>} & \texttt{Box<T>} & no null state; \texttt{Option<Box<T>{}>{}} if you need one\\
\texttt{std::shared\_ptr<T>} & \texttt{Rc<T>} / \texttt{Arc<T>} & \texttt{Rc} is single-thread only; compiler enforces it\\
\texttt{T*} (nullable, non-owning) & \texttt{Option<\&T>} & no null deref possible\\
\texttt{std::optional<T>} & \texttt{Option<T>} & \texttt{Option} must be opened; no implicit bool\\
\texttt{std::variant<A, B>} & \texttt{enum} with data & exhaustive \texttt{match}, no \texttt{std::visit} ceremony\\
\texttt{std::map<K, V>} & \texttt{BTreeMap<K, V>} & needs \texttt{K: Ord} — hence the whole Step 4\\
\texttt{std::unordered\_map<K, V>} & \texttt{HashMap<K, V>} & needs \texttt{K: Hash + Eq}; not sorted\\
\texttt{std::function<R(A)>} & \texttt{Box<dyn Fn(A) -> R>} & closures are types; \texttt{Fn} tier is explicit\\
abstract base + \texttt{virtual} & \texttt{trait} + \texttt{dyn Trait} & no inheritance of \textbf{data}, only of contract\\
template \texttt{<typename T>} & generic \texttt{<T: Bound>} & bounds checked at definition, not instantiation\\
\texttt{const\&} member fn & \texttt{\&self} & borrow checker enforces it, not convention\\
\texttt{mutable} member & \texttt{Cell} / \texttt{RefCell} & opt-in and visible in the type (Appendix C)\\
\texttt{throw} / \texttt{try} & \texttt{Result} + \texttt{?} & no unwinding across FFI; errors are return values\\
\texttt{std::move(x)} & just \texttt{x} & move is the default; no moved-from husk to read\\
RAII destructor & \texttt{Drop} & same idea; runs at last use, not scope end (NLL)\\
\texttt{assert} & \texttt{assert!} / \texttt{debug\_assert!} & \texttt{debug\_assert!} compiles out in release\\
\end{tabular}
\end{center}
\subsection{From Python}
\label{sec:orgabc436f}

\begin{center}
\begin{tabular}{lll}
Python & Rust & The difference that matters\\
\hline
\texttt{list} & \texttt{Vec<T>} & one element type, no mixing\\
\texttt{dict} & \texttt{HashMap} / \texttt{BTreeMap} & key needs \texttt{Hash+Eq} or \texttt{Ord}\\
\texttt{tuple} & \texttt{(A, B)} & fixed arity in the type\\
\texttt{None} & \texttt{Option::None} & a different \textbf{type}, not a value of every type\\
\texttt{raise} / \texttt{except} & \texttt{Result} + \texttt{?} & callers see failure in the signature\\
duck typing & \texttt{trait} bounds & the contract is written down and checked\\
\texttt{@dataclass} & \texttt{struct} + \texttt{\#[derive(...)]} & derives are opt-in per capability\\
list comprehension & \texttt{.iter().map().filter().collect()} & lazy until \texttt{collect}\\
\texttt{for i in range(len(x))} & \texttt{for item in \&x} & idiomatic Rust rarely indexes (\texttt{perf-iter-over-index})\\
\texttt{copy.deepcopy} & \texttt{.clone()} & explicit, visible, and usually avoidable\\
GIL / ``just use threads'' & \texttt{Send} / \texttt{Sync} & data races are compile errors (Appendix B)\\
\texttt{isinstance} dispatch & \texttt{match} on an \texttt{enum} & exhaustive; new variant = compile error at every site\\
\end{tabular}
\end{center}

The Python row worth staring at is the last one. In Python, adding a case means
hunting for every \texttt{isinstance} chain. In Rust, adding an enum variant makes
every non-wildcard \texttt{match} fail to compile until you handle it. That is the
single biggest reason this doc keeps pushing \texttt{enum} over \texttt{u8} — it's not
tidiness, it's a refactoring superpower.
\section{Appendix H — debugging the FFI seam}
\label{sec:orgb59663d}

Everything above assumes things work. This is the appendix for 3pm on the day
the linker starts shouting. Rust's compile errors are excellent; cxx's and the
linker's are not, so pattern-matching on the symptom is the fastest route.
\subsection{First: look at what cxx generated}
\label{sec:org5eb6b3a}

Two commands answer most questions:

\begin{minted}[]{sh}
cargo expand --lib ffi        # the Rust side the bridge macro wrote
cat target/*/cxxbridge/direct-rs/src/lib.rs.h   # the C++ side it wrote
\end{minted}

Every ``is it \texttt{rust::Box<T>} or \texttt{rust::box<T>}, does the slice come in as
\texttt{rust::Slice<const Pose>}'' question in this doc is answered by that header, for
\textbf{your} cxx version, in seconds. This is the smoke-crate advice from ``Review
policy'', but you don't need a smoke crate — you need to read the artifact.
\subsection{Symptom → cause table}
\label{sec:org24fba6e}

\begin{center}
\begin{tabular}{lll}
Symptom & Cause & Fix\\
\hline
\texttt{undefined reference to direct\_rs::new\_direct\_optimizer} & staticlib not linked, or name mangling differs & check \texttt{crate-type}{[}``staticlib'']= + \texttt{target\_link\_libraries}\\
\texttt{undefined reference to CppCost::evaluate6} & declared in bridge, not defined in C++ & \texttt{evaluate6} must be a real (non-inline-only) member\\
cxx: ``unsupported type'' on a bridge signature & tried to pass a non-shared, non-opaque type & shared struct of PODs, or opaque + reference; nothing else\\
cxx: ``self must be a reference'' & \texttt{fn run(self, ...)} on an opaque type & \texttt{\&self} / \texttt{\&mut self} only (see Ownership/threading)\\
cxx: function pointer / callback rejected & C++→Rust fn pointers unimplemented & invert to the \texttt{ProgressSink} shape (M2)\\
process dies with \texttt{thread panicked ... note: aborting} & a panic reached an \texttt{extern} boundary & \texttt{?} the error, return \texttt{ok:false}; never \texttt{unwrap} at the edge\\
C++ sees garbage in a \texttt{Pose} field & field \textbf{order} differs between the two decls & shared struct means ONE decl — don't hand-write the C++ side\\
optimum matches, \texttt{calls} doesn't & budget truncation differs mid-batch & see Appendix E, item 2 — settle the open question\\
\texttt{Pose} values fine, costs all \texttt{inf} & handle outlived / GPU state not set up & the \texttt{CppCost} lifetime rule; check it's alive across \texttt{run}\\
works in debug, wrong in release & UB, or an \texttt{f64} comparison/=debug\textsubscript{assert}!= & run under sanitizers; don't put logic in \texttt{debug\_assert!}\\
\end{tabular}
\end{center}
\subsection{Tools worth having configured before you need them}
\label{sec:orgc4672a2}

\begin{itemize}
\item \texttt{RUST\_BACKTRACE=1} — set it in the launch environment permanently. A Rust
panic without a backtrace tells you almost nothing.
\item \texttt{nm -C libdirect\_rs.a | grep direct\_rs} — proves the symbol exists before you
blame CMake.
\item \texttt{gdb} / \texttt{lldb} — both step across the boundary; Rust frames show up with
mangled-but-readable names. Breaking on \texttt{rust\_panic} catches the abort case
\textbf{at the panic}, not after.
\item \texttt{RUSTFLAGS}``-Z sanitizer=address''= (nightly) — worth one run once the batch
slice path is live, because that's the only place raw memory is shared.
\item \texttt{cargo tree} — when two crates disagree on a \texttt{cxx} version, which produces
spectacularly confusing layout mismatches.
\end{itemize}
\subsection{The panic firewall (do this once, in \texttt{lib.rs})}
\label{sec:orgbde4324}

Because a panic crossing FFI is undefined behavior and the design says ``never
let Rust panic across FFI'', make that structural rather than aspirational:

\begin{minted}[]{rust}
fn run(&mut self, cost: &CppCost) -> ffi::RunOutcome {
    std::panic::catch_unwind(std::panic::AssertUnwindSafe(|| {
        self.run_inner(cost)
    }))
    .unwrap_or_else(|_| ffi::RunOutcome { ok: false, ..Default::default() })
}
\end{minted}

Caveats worth understanding rather than copying blindly: \texttt{catch\_unwind} does
nothing under \texttt{panic=abort}, and \texttt{AssertUnwindSafe} is you promising the state
isn't left half-mutated. It's a net, not a design. The design is \texttt{Result}
everywhere inside (Appendix D) — this catches the bug you didn't predict, on a
customer's machine, instead of aborting the whole application.

\begin{itemize}
\item \textbf{\textbf{Book}}: ch 09-01 (panic and unwinding), ch 19-01 (=unsafe=/FFI).
\item \textbf{\textbf{rust-skills}}: \texttt{err-no-unwrap-prod}, \texttt{unsafe-safety-comment}.
\end{itemize}

Stop. \textbf{\textbf{Check}} — the C++ side reports a plausible optimum but \texttt{out.ok =} false=
every run. Where do you look first? (The \texttt{Result=→=ok} flatten in \texttt{run\_ffi}:
something inside is returning \texttt{Err} \textbf{after} best was already updated — which is
exactly the information the flatten throws away. Log the \texttt{OptimizerError}
on the Rust side before you drop it; this is the ``what is lost at the edge''
from M8, biting.)
\section{Rust idiom glossary + Brown Book index}
\label{sec:org5a07f52}

A cross-reference table: the concept this FFI handshake touches, where to read
it in the Brown Book, the rust-skills rule, and a one-line ``normie → idiomatic''
teaser. Use it to jump straight to the depth you need.

\begin{center}
\begin{tabular}{llll}
Rust concept & Brown Book & rust-skills rule & normie → idiomatic guardrail\\
\hline
ownership / move semantics & ch 04-01 & own-borrow-over-clone & ``don't copy, don't pass by value — move or borrow''\\
borrowing = \&T / \&mut T & ch 04-02 & own-borrow-over-clone & ``\&mut is exclusive; \& is shared''\\
slices = \&[T] (fat pointer) & ch 04-03 & own-slice-over-vec & ``accept \&[T], not \&Vec<T>''\\
structs (value types) & ch 05 & own-copy-small & ``small POD-alleys by value; big/owned on heap''\\
enums + exhaustive match & ch 06 & type-enum-states & ``use enum for mutually-exclusive states, not u8''\\
Result / ? / no exceptions & ch 09 & err-result-over-panic & ``errors are values; caller decides''\\
Box<T> ≈ unique\textsubscript{ptr} & ch 15-01 & own-move-large & ``only box when you need heap or a stable address''\\
traits as contracts & ch 10-02 & trait-dyn-vs-generic & ``\texttt{Cost} is a trait, not a base class''\\
dyn Trait / trait objects & ch 17 & trait-object-safety & ``use dyn only when runtime dispatch is needed''\\
closures: Fn/FnMut/FnOnce & ch 13 & closure-fn-trait-bounds & ``request the least-restrictive Fn you need''\\
iterators (lazy, chainable) & ch 13 & perf-iter-over-index & ``map/filter/collect, not index loops''\\
unsafe + FFI / \#[repr(C)] & ch 19 & unsafe-safety-comment & ``boundary-only; never leak a panic''\\
error type with thiserror & ch 09-02 & err-thiserror-lib & ``derive Error; use ? everywhere inside''\\
panic = abort across FFI & ch 09-01 & err-no-unwrap-prod & ``unwrap only on verified invariants''\\
newtype for IDs/units & ch 19-03? & type-newtype-ids & ``BoxBudget(u32), not bare u32''\\
Default derive for options & ch 05 & api-default-impl & ``Default == C++ default ctor''\\
\end{tabular}
\end{center}

Key Brown Book sections to open (each is a chapter page; substitute the number
into \url{https://rust-book.cs.brown.edu/chNN-0X-...}.html):

\begin{enumerate}
\item \texttt{ch04-02-references-and-borrowing.html} — the borrow table (move/\&T/\&mut T).
\item \texttt{ch09-02-recoverable-errors-with-result.html} — Result + ?.
\item \texttt{ch10-02-traits-defining-shared-behavior.html} — the \texttt{Cost} trait.
\item \texttt{ch13-01-closures.html} — Fn traits + capture.
\item \texttt{ch15-01-box-point-to-data-on-the-heap.html} — why Box here.
\item \texttt{ch19-01-unsafe-rust.html} — what \texttt{unsafe extern} gives up.
\end{enumerate}

How to grow with this doc: whenever you implement a piece and hit a borrow
checker error, look it up in the glossary's ``normie → idiomatic'' column first;
the Brown chapter it points to walks through the visual/physical model, and the
rust-skills rule gives the sharper 1-3 sentence rule. Re-read ch 4 after every
few-up milestone — ownership is the book's spine, and the FFI handshake above
is one long exercise in it.
\subsection{Borrow-checker error lookup (keep this beside the Brown ch 4 quizzes)}
\label{sec:org0e75856}

The fastest way to grow from each compiler rejection is to map the error to a
cause and a move. This is the same canon as Appendix A, but organized as a
diagnostic cheat-sheet with the \emph{normie} reflex you'll have and the \emph{idiomatic}
response.

\begin{center}
\begin{tabular}{llll}
Error code & You (normie) wrote & What Rust means & The idiomatic move\\
\hline
E0382 & reuse a \texttt{Hyperbox} after \texttt{pop\_first()} & ``use of moved value'' — you gave it & reorder so the last use comes first, or \texttt{.clone()} only the small bits\\
 & moved it & to the map / to trisect & you still need\\
E0502 & \texttt{\&self} borrow alive while you call =\&mut & "cannot borrow as mutable; it's & scope the shared borrow (block or \texttt{drop}) before the \texttt{\&mut}; or split\\
 & self=-mutating method & also borrowed as immutable" & into a \&=self= gather + \&=mut= apply\\
E0499 & two live \texttt{\&mut} into the tree (nested \texttt{get\_mut} loop) & ``cannot borrow as mutable more than once'' & restructure so you never hold two \texttt{\&mut} into the same map\\
E0505 & "cannot move out of \texttt{*x} behind a =\&\\
\end{tabular}
\end{center}
reference" inside the gather                 | you can't move out of a shared
borrow                          | copy the small fields (\texttt{*size\_key}.into\textsubscript{inner}()=), not the whole box   |
\begin{center}
\begin{tabular}{llll}
E0596 & calling a mutating method through \texttt{\&self} & ``cannot borrow as mutable'' (receiver is \texttt{\&self}) & change the method to \texttt{\&mut self}, or return data instead of scribbling\\
E0753 & FFI panic crossed the boundary & ``panic in a function that cannot unwind'' & never let a panic reach an \texttt{extern}; catch it, return a status struct\\
\end{tabular}
\end{center}

Read the first column as ``the code are paths where the borrow checker caught a
C++-shaped habit.'' Each fix is the corresponding Rust-shaped habit.
\subsection{Self-check index (the Roadmap's questions, collected)}
\label{sec:orge893455}

Answers you should be able to give — by milestone — without looking. If one is
fuzzy, that's the part to re-read before implementing.

\begin{center}
\begin{tabular}{lll}
\# & Milestone & The self-check question\\
\hline
M0 & split & who owns cost / search / thread / UI / budget?\\
M1 & types & why \texttt{Pose} is \texttt{Copy} but a \texttt{Vec}-wrapping type isn't?\\
M2 & BTreeMap & why can't \texttt{f64} be a key; what's the minimal fix?\\
M3 & gather & why \texttt{\&self}; what breaks if you clone the box out?\\
M4 & POH & how do you know \texttt{convex\_hull} didn't touch the tree?\\
M5 & trisect & why pop/remove OUT before trisecting? why no in-place map mutation?\\
M6 & Cost & why a trait? when \texttt{impl Fn} vs \texttt{Box<dyn Cost>}?\\
M7 & cxx & what two things can't cross a cxx boundary?\\
M8 & errors & Result-in vs status-out: what's lost at the edge and why OK?\\
M9 & assemble & walk one \texttt{Hyperbox} borrowing/move end-to-end.\\
A & borrow & identify E0505 vs E0502 vs E0499 when you see them.\\
B & lifetime & why \texttt{CppCost} isn't auto-\texttt{Send/Sync}; how one thread dodges it.\\
C & RefCell & run the 3-step checklist before adding \texttt{RefCell} for the counter.\\
D & Option & three correct ways to get an \texttt{f64} out of an \texttt{Option<f64>}, and when each.\\
E & run loop & trace calls/FFI-hops/callback order for one iteration of a 4-box hull.\\
F & testing & which test tier catches a tie-break divergence, and why not the others?\\
G & translate & \texttt{std::function} → what? \texttt{std::variant} → what? what changes, not just the name?\\
H & debugging & first two commands when the linker or the layout goes wrong.\\
\end{tabular}
\end{center}

Use this table as a study checklist: each row is \textasciitilde{}10 minutes of reading in the
linked chapter plus one quiz. When you can answer all eighteen, you own every
concept the handshake runs on.

And the honest ordering, if you only have a week: Toolchain → exercises E1-E7
→ Appendix D → Appendix A → Appendix E → Appendix F Tier 3. That path ends with
a Rust DIRECT that provably matches the C++ on a recorded cost sequence, which
is the only milestone that actually de-risks the port. The rest of this doc is
reference you'll re-read as the compiler sends you back to it.
\end{document}
